% !TEX root = ../../main.tex
%\setchapterimage{figure/chapter5封面图.jpg}
\setchapterintro{}
\chapter{线性变换以及相似理论}

\section{不变子空间}

\begin{definition}
	设 $T\in\End(V)$，$W$ 是 $V$ 的子空间.若对任意 $\alpha\in W$，都有
	\[
	T(\alpha)\in W,
	\]
	则称 $W$ 是 $T$ 的一个不变子空间，或称 $W$ 是 $V$ 的一个 $T$-子空间.
	
	若 $W$ 是一个 $T$-子空间，则有线性变换 $T|_W\in\End(W)$：
	\[
	W\ni\alpha\longmapsto T(\alpha)\in W.
	\]
	称 $T|_W$ 为 $T$ 在 $W$ 上的限制.
\end{definition}

对于 $T\in\End(V)$，显然地 $0,V$ 是 $T$-子空间，称为 $T$ 的平凡不变子空间.
注意：一个线性映射可能没有任何非平凡不变子空间，如 $\mathbb{R}^2$ 上的旋转.

\begin{proposition}
	设 $T\in\End(V)$，则 $\ker T$ 与 $\Im T$ 是 $T$-子空间.
\end{proposition}

\begin{proof}
	\begin{enumerate}
	\item 对任意 $\alpha\in\ker T$，有
	\[
	T\alpha=0\in\ker T,
	\]
	故 $T|_{\ker T}=0$.
	
	\item 对任意 $\beta\in\Im T$，有
	\[
	T(\beta)\in\Im T,
	\]
	故 $T|_{\Im T}=T$.
\end{enumerate}
\end{proof}


\begin{corollary}
设 $T\in\End(V)$，$f\in\F[x]$，则 $\ker f(T)$ 与 $\Im f(T)$ 是 $T$-子空间.
\end{corollary}

\begin{example}
设 $T\in\End(V)$.若 $W_1,W_2$ 是 $T$-子空间，则 $W_1\cap W_2$ 与 $W_1+W_2$ 仍是 $T$-子空间.
\end{example}

\begin{example}
设 $T\in\End(V)$，$\alpha\in V$，则
\[
\langle \alpha,T\alpha,T^2\alpha,\cdots,T^n\alpha,\cdots\rangle
\]
是一个 $T$-子空间，称为 Krylov 空间，或循环子空间.
\end{example}

\begin{example}
设 $S,T\in\End(V)$ 且 $ST=TS$，则 $\ker S$ 与 $\Im S$ 是 $T$-子空间.
\end{example}

\begin{proof}
对任意 $\alpha\in\ker S$，有
\[
S(T\alpha)=TS(\alpha)=0,
\]
故 $T\alpha\in\ker S$.对任意 $\beta\in\Im S$，存在 $\alpha\in V$ 使得 $S(\alpha)=\beta$，于是
\[
T\beta=TS(\alpha)=ST(\alpha)\in\Im S.
\]
\end{proof}

\begin{example}
设 $T\in\End(V)$，$f\in\F[x]$，令
\[
S=T|_{\ker f(T)}.
\]
问 $S$ 的零化多项式.
\end{example}

\begin{proof}
对任意 $\alpha\in\ker f(T)$，有
\[
f(T)(\alpha)=0.
\]
于是
\[
f(S)(\alpha)=f\bigl(T|_{\ker f(T)}\bigr)(\alpha)=f(T)(\alpha)=0.
\]
故 $f(S)=0$.
\end{proof}

\begin{example}
设 $f=f_1\cdots f_s$，其中 $f_1,\cdots,f_s$ 两两互素.若 $f(T)=0$，则
\[
V=\ker f_1(T)\oplus\cdots\oplus\ker f_s(T).
\]
因此 $V$ 是 $s$ 个不变子空间的直和.
\end{example}

\begin{proposition}
设 $T\in\End(V)$，且
\[
V=V_1\oplus\cdots\oplus V_s,
\]
其中 $V_i$ 是 $T$-子空间.设 $T|_{V_i}$ 的极小多项式为 $m_i(x)$，则 $T$ 的极小多项式为
\[
[m_1(x),\cdots,m_s(x)].
\]
\end{proposition}

\begin{proof}
令
\[
f(x)=[m_1(x),\cdots,m_s(x)],
\qquad
f(x)=m_i(x)h_i(x),\quad 1\leq i\leq s.
\]
对任意 $\beta_i\in V_i$，都有
\[
\begin{aligned}
	f(T)(\beta_i)
	&=h_i(T)m_i(T)(\beta_i)\\
	&=h_i(T)m_i\bigl(T|_{V_i}\bigr)(\beta_i)=0.
\end{aligned}
\]
对任意 $\alpha\in V$，存在唯一的 $\alpha_i\in V_i$，使得
\[
\alpha=\alpha_1+\cdots+\alpha_s.
\]
于是
\[
f(T)(\alpha)=f(T)(\alpha_1)+\cdots+f(T)(\alpha_s)=0.
\]
故 $f$ 是 $T$ 的零化多项式.

设 $m(x)$ 是 $T$ 的极小多项式.由 $f(T)=0$，得 $m\mid f$.又对任意 $\gamma_i\in V_i$，有
\[
0=m(T)(\gamma_i)=m\bigl(T|_{V_i}\bigr)(\gamma_i),
\]
故 $m_i(x)\mid m(x)$，$1\leq i\leq s$，从而
\[
[m_1(x),\cdots,m_s(x)]\mid m(x).
\]
因此 $m=f$.
\end{proof}
\section{特征子空间、特征值、特征向量}

先看一个最基本的例子.设 $T\in\End(V)$，$a\in\F$，令
\[
V_a=\{\gamma\in V:T\gamma=a\gamma\}.
\]
显然
\[
V_a=\ker(T-a\id_V),
\]
故 $V_a$ 是 $V$ 的一个子空间.并且 $V_a$ 是 $T$-子空间，因为若 $\gamma\in V_a$，则
\[
T(T\gamma)=T(a\gamma)=aT\gamma,
\]
所以 $T\gamma\in V_a$.同时
\[
T|_{V_a}=a\id_{V_a},
\]
即 $T$ 在 $V_a$ 上就是数乘变换.

\begin{definition}
	设 $T\in\End(V)$，$a\in\F$.若
	\[
	V_a=\ker(T-a\id_V)\neq0,
	\]
	则称 $a$ 为 $T$ 的一个特征值，称 $V_a$ 为 $T$ 属于特征值 $a$ 的特征子空间.$V_a$ 中的非零向量称为 $T$ 属于特征值 $a$ 的特征向量.
	
	若 $\dim V_a<\infty$，则称
	\[
	\dim V_a
	\]
	为特征值 $a$ 的几何重数.
\end{definition}

注意：$0$ 不称为特征向量.严格说，属于 $a$ 的特征向量全体为 $V_a\setminus\{0\}$.当 $\F$ 为无限域时，一个特征值若存在，则有无穷多个特征向量.

另外一个特征向量只能属于一个特征值

	设 $T\in\End(V)$，$a,b\in\F$，且 $a\neq b$.若 $\xi$ 同时属于特征值 $a,b$

则$T\xi=a\xi$ 且 $T\xi=b\xi$，则
	\[
	(a-b)\xi=0.
	\]
	由 $a-b\neq0$，得 $\xi=0$.

从而
\[
V_a\cap V_b=0.\text{且}V_a + V_b \text{是直和}
\]

\begin{proposition}
	设 $T\in\End(V)$，$a_1,\cdots,a_s\in\F$ 是 $T$ 的互不相同的特征值，则
	\[
	V_{a_1}+\cdots+V_{a_s}
	\]
	是直和.
\end{proposition}

\begin{proof}
	设
	\[
	0=\xi_1+\cdots+\xi_s,\qquad \xi_i\in V_{a_i}.
	\]
	对上式依次作用 $T^k$，$k=0,1,\cdots,s-1$，得
	\[
	\sum_{i=1}^s a_i^k\xi_i=0,
	\qquad k=0,1,\cdots,s-1.
	\]
	即
	\[
	\begin{pmatrix}
		1&1&\cdots&1\\
		a_1&a_2&\cdots&a_s\\
		\vdots&\vdots&&\vdots\\
		a_1^{s-1}&a_2^{s-1}&\cdots&a_s^{s-1}
	\end{pmatrix}
	\begin{pmatrix}
		\xi_1\\ \xi_2\\ \vdots\\ \xi_s
	\end{pmatrix}
	=0.
	\]
	左边的 Vandermonde 矩阵行列式为
	\[
	\prod_{1\leq i<j\leq s}(a_j-a_i)\neq0,
	\]
	故 $\xi_1=\cdots=\xi_s=0$.因此和为直和.
\end{proof}


\begin{definition}
	若 $V$ 可以分解为 $T$ 的特征子空间的直和，即存在互不相同的特征值 $a_1,\cdots,a_s$，使得
	\[
	V=V_{a_1}\oplus\cdots\oplus V_{a_s},
	\]
	则称 $T$ 可以对角化，或称 $T$ 是可对角化的.
\end{definition}
后面我们将证明，可对角化等价于可以取到一组基，使得 $T$ 的矩阵为对角矩阵.也即：

若 $\xi_{i1},\cdots,\xi_{it_i}$ 是 $V_{a_i}$ 的一组基，则
\[
\xi_{11},\cdots,\xi_{1t_1},\xi_{21},\cdots,\xi_{2t_2},\cdots,
\xi_{s1},\cdots,\xi_{st_s}
\]
是 $V$ 的一组基，且 $T$ 在这组基下的矩阵为
\[
\begin{pmatrix}
	a_1I_{t_1}&&&0\\
	&a_2I_{t_2}&&\\
	&&\ddots&\\
	0&&&a_sI_{t_s}
\end{pmatrix}.
\]





\begin{example}
	设 $f\in\F[x]$，定义
	\[
	T:\F[x]\longrightarrow\F[x],\qquad g\longmapsto fg.
	\]
	若 $f$ 不是常数多项式，则 $T$ 没有特征值.
\end{example}

\begin{proof}
	若 $Tg=ag$，则
	\[
	fg=ag\rightarrow f = a
	\]
	
\end{proof}

\begin{example}
	设
	\[
	T\binom{x}{y}=\binom{y}{-x}.
	\]
	在 $\mathbb{R}^2$ 上，$T$ 没有特征值；在 $\mathbb{C}^2$ 上，$T$ 有特征值 $i,-i$，并且
	\[
	V_i=\left\langle\binom{1}{i}\right\rangle,
	\qquad
	V_{-i}=\left\langle\binom{1}{-i}\right\rangle.
	\]
\end{example}

\begin{proof}
	若在 $\mathbb{R}^2$ 中有
	\[
	\binom{y}{-x}=a\binom{x}{y},
	\]
	则
	\[
	\begin{cases}
		y=ax,\\
		-x=ay.
	\end{cases}
	\]
	故 $(a^2+1)x=0$，进而 $x=y=0$，矛盾.因此在 $\mathbb{R}^2$ 中无特征值.在 $\mathbb{C}^2$ 中直接代入可得 $i,-i$.
\end{proof}

\begin{example}
	令
	\[
	T\in\End(M_n(\F)),
	\qquad T(A)=A'.
	\]
	则 $T$ 的特征值为 $1,-1$，且
	\[
	V_1=S_n=\{A\in M_n(\F):A'=A\},
	\]
	\[
	V_{-1}=A_n=\{A\in M_n(\F):A'=-A\}.
	\]
	并且
	\[
	M_n(\F)=S_n\oplus A_n.
	\]
\end{example}
\begin{corollary}
	设 $T\in\End(V)$，$a_1,\cdots,a_s\in\F$ 是 $T$ 的互不相同的特征值
	, $\xi_{i1},\cdots,\xi_{it_i}$ 是 $V_{a_i}$ 的一组基，则
	\[
	\xi_{11},\cdots,\xi_{1t_1},\xi_{21},\cdots,\xi_{2t_2},\cdots,
	\xi_{s1},\cdots,\xi_{st_s}
	\]
	是 $V=V_{a_1}\oplus\cdots\oplus V_{a_s}$ 的一组基
\end{corollary}
\begin{corollary}
	有限维线性空间上的任意线性变换至多有有限个特征值
\end{corollary}
无限维情形中，特征值可以有无穷多个.例如在 $C^\infty(a,b)$ 上令
\[
D=\frac{d}{dx},
\]
则对任意 $\lambda\in\F$，都有
\[
D(e^{\lambda x})=\lambda e^{\lambda x},
\]
故 $\lambda$ 是 $D$ 的特征值.

\begin{lemma}
	设 $T\in\End(V)$，$a$ 是 $T$ 的一个特征值，$\xi\in V_a$，$f(x)\in\F[x]$.则
	\[
	f(T)\xi=f(a)\xi.
	\]
	特别地，若 $\xi\neq0$，则 $\xi$ 也是 $f(T)$ 属于特征值 $f(a)$ 的特征向量.
\end{lemma}

\begin{proof}
	由 $T\xi=a\xi$ 可得
	\[
	T^k\xi=a^k\xi,
	\qquad k\geq0.
	\]
	若 $f(x)=c_0+c_1x+\cdots+c_rx^r$，则
	\[
	f(T)\xi=\sum_{k=0}^r c_kT^k\xi
	=\sum_{k=0}^r c_ka^k\xi=f(a)\xi.
	\]
\end{proof}

反过来，$f(T)$ 的特征向量不一定都是 $T$ 的特征向量.

\begin{example}
	在 $C^\infty(a,b)$ 上令
	\[
	L=\frac{d^2}{dx^2}+p\frac{d}{dx}+qI,
	\qquad p,q\in\mathbb{R}.
	\]
	对任意 $\lambda\in\mathbb{R}$，有
	\[
	L(e^{\lambda x})=(\lambda^2+p\lambda+q)e^{\lambda x}.
	\]
	因此 $e^{\lambda x}$ 是 $L$ 的特征向量，对应特征值为
	\[
	\lambda^2+p\lambda+q.
	\]
\end{example}

\begin{example}
	令
	\[
	T:M_n(\F)\longrightarrow M_n(\F),
	\qquad T(A)=A'-aA.
	\]
	记 $T_0(A)=A'$，则 $T=T_0-aid
	$.由上引理可知，$T$ 的特征值为
	\[
	1-a,
	\qquad -1-a,
	\]
	其特征子空间分别为 $S_n$ 与 $A_n$.
\end{example}

\begin{lemma}
	设 $T\in\End(V)$，$a_1,\cdots,a_s$ 是 $T$ 的互不相同的特征值.若 $W$ 是 $V$ 的一个 $T$-子空间，则
	\[
	W\cap(V_{a_1}\oplus\cdots\oplus V_{a_s})
	=(W\cap V_{a_1})\oplus\cdots\oplus(W\cap V_{a_s}).
	\]
\end{lemma}

\begin{proof}
	只需证左边包含于右边.任取
	\[
	\alpha\in W\cap(V_{a_1}\oplus\cdots\oplus V_{a_s}),
	\qquad
	\alpha=\alpha_1+\cdots+\alpha_s,
	\quad \alpha_i\in V_{a_i}.
	\]
	因 $W$ 是 $T$-子空间，$T^k\alpha\in W$，$k=0,1,\cdots,s-1$.又
	\[
	\begin{pmatrix}
		\alpha\\ T\alpha\\ \vdots\\ T^{s-1}\alpha
	\end{pmatrix}
	=
	\begin{pmatrix}
		1&1&\cdots&1\\
		a_1&a_2&\cdots&a_s\\
		\vdots&\vdots&&\vdots\\
		a_1^{s-1}&a_2^{s-1}&\cdots&a_s^{s-1}
	\end{pmatrix}
	\begin{pmatrix}
		\alpha_1\\ \alpha_2\\ \vdots\\ \alpha_s
	\end{pmatrix}.
	\]
	Vandermonde 矩阵可逆，故每个 $\alpha_i$ 都是 $\alpha,T\alpha,\cdots,T^{s-1}\alpha$ 的线性组合，从而 $\alpha_i\in W$.于是 $\alpha_i\in W\cap V_{a_i}$.
\end{proof}

\begin{corollary}
	若 $T\in\End(V)$ 可对角化，且 $W$ 是 $V$ 的一个 $T$-子空间，则 $T|_W$ 也可对角化.
\end{corollary}

\begin{proof}
	设
	\[
	V=V_{a_1}\oplus\cdots\oplus V_{a_s}.
	\]
	由上一个引理，
	\[
	W=W\cap V=(W\cap V_{a_1})\oplus\cdots\oplus(W\cap V_{a_s}).
	\]
	而每个非零的 $W\cap V_{a_i}$ 都是 $T|_W$ 的特征子空间，故 $T|_W$ 可对角化.
\end{proof}
\begin{example}
	设 $T \in \End(V)$ 可对角化，且 $\dim V = n$，且 $T$ 的每个特征子空间都是1维的.求 $T$ 的全部不变子空间的个数.
	
	
\end{example}
\begin{solution}
	\begin{align*}
		& \text{有 } n \text{ 个特征子空间}. \\
		& \text{任取 } k \text{ 个求和 } \mathrm{C}_n^k \Rightarrow \sum_{k=0}^n \mathrm{C}_n^k = 2^n. \\
		& \text{再证没有别的不变子空间}. \\
		& \forall W \text{，证明 } W \text{ 可由特征子空间和出：} \\
		W &= W \cap V = W \cap (V_{\lambda_1} \oplus \cdots \oplus V_{\lambda_n}) \\
		&= W \cap V_{\lambda_1} \oplus \cdots \oplus W \cap V_{\lambda_n}
	\end{align*}
	故 $W$ 是这些交的直和，因此只能是若干特征子空间的直和，个数为 $2^n$.
\end{solution}


\begin{example}
	设 $S,T\in\End(V)$ 且 $ST=TS$.若 $a$ 是 $S$ 的特征值，则 $S$ 的特征子空间
	\[
	V_a(S)=\ker(S-aI)
	\]
	是 $T$-子空间.
\end{example}

\begin{proof}
	任取 $\xi\in V_a(S)$，则
	\[
	S(T\xi)=T(S\xi)=T(a\xi)=aT\xi.
	\]
	故 $T\xi\in V_a(S)$.
\end{proof}

\section{方阵特征多项式与三角化}

这一节讨论方阵特征值、特征向量、特征多项式以及不变子空间和矩阵形状之间的关系.

\begin{definition}
	设 $A\in M_n(\F)$.称
	\[
	f_A(\lambda)=\det(\lambda E_n-A)
	\]
	为 $A$ 的特征多项式.若 $a\in\F$，则
	\[
	a\text{ 是 }A\text{ 的特征值}
	\Longleftrightarrow
	\mathcal N(A-aE_n)\neq0
	\Longleftrightarrow
	\det(aE_n-A)=0.
	\]
	特征值 $a$ 作为 $f_A(\lambda)$ 的根的重数称为 $a$ 的代数重数.
\end{definition}

若 $A=(a_{ij})$，则
\[
f_A(\lambda)=
\begin{vmatrix}
	\lambda-a_{11}&-a_{12}&\cdots&-a_{1n}\\
	-a_{21}&\lambda-a_{22}&\cdots&-a_{2n}\\
	\vdots&\vdots&\ddots&\vdots\\
	-a_{n1}&-a_{n2}&\cdots&\lambda-a_{nn}
\end{vmatrix}.
\]
因此 $f_A(\lambda)$ 是一个关于 $\lambda$ 的 $n$ 次首一多项式，并且
\[
f_A(\lambda)=\lambda^n-\operatorname{tr}(A)\lambda^{n-1}+\cdots+(-1)^n\det A.
\]

若 $\F$ 是代数闭域，则 $f_A$ 在 $\F$ 中必有根，所以每个 $A\in M_n(\F)$ 都至少有一个特征值.但在 $\mathbb R$ 上不一定成立.

\begin{example}
	设
	\[
	A=\begin{pmatrix}0&-1\\ 1&0\end{pmatrix}.
	\]
	则
	\[
	f_A(\lambda)=\lambda^2+1.
	\]
	在 $\mathbb R$ 上，$A$ 没有特征值；在 $\mathbb C$ 上，$A$ 的特征值为 $i,-i$.
\end{example}

\begin{example}
	若 $A$ 是幂零矩阵，即存在 $m>0$ 使得 $A^m=0$，则 $A$ 唯一可能的特征值是 $0$.
\end{example}

\begin{proof}
	若 $A\xi=a\xi$ 且 $\xi\neq0$，则
	\[
	0=A^m\xi=a^m\xi.
	\]
	故 $a^m=0$，从而 $a=0$.
\end{proof}

\begin{example}
	设 $A\in M_n(\F)$.则
	\[
	A\in \GL_n(\F)
	\Longleftrightarrow
	0\text{ 不是 }A\text{ 的特征值}.
	\]
	更一般地，$a$ 是 $A$ 的特征值当且仅当 $A-aE_n$ 不可逆.
\end{example}
	


\begin{corollary}
	若 $A\in \GL_n(\F)$，则 $a$ 是 $A$ 的特征值当且仅当 $a\neq0$ 且 $a^{-1}$ 是 $A^{-1}$ 的特征值.
\end{corollary}

\begin{proof}
	若 $A\xi=a\xi$，$\xi\neq0$，则 $a\neq0$，且
	\[
	A^{-1}\xi=a^{-1}\xi.
	\]
	反过来同理.
\end{proof}
\begin{example}
	若 $A,B\in M_n(\F)$，则 $AB$ 与 $BA$ 有相同的非零特征值.
\end{example}

\begin{proof}
	由\[|E_n - \frac{1}{\lambda}AB| =|E_n - B\frac{1}{\lambda}A|\]易知
\end{proof}


\begin{example}
	若 $A$ 是上三角矩阵
	\[
	A=\begin{pmatrix}
		a_{11}&*&\cdots&*\\
		0&a_{22}&\cdots&*\\
		\vdots&\vdots&\ddots&\vdots\\
		0&0&\cdots&a_{nn}
	\end{pmatrix},
	\]
	则
	\[
	f_A(\lambda)=(\lambda-a_{11})(\lambda-a_{22})\cdots(\lambda-a_{nn}).
	\]
	故上三角矩阵的特征值就是对角线上的元素.
	
	若是准上三角阵
	\[
	A=\begin{pmatrix}
		A_{11}&*&\cdots&*\\
		0&A_{22}&\cdots&*\\
		\vdots&\vdots&\ddots&\vdots\\
		0&0&\cdots&A_{nn}
	\end{pmatrix},
	\]
	则\[f(\lambda) = f_{A_{11}}(\lambda)\cdots f_{A_{nn}}(\lambda)\]
\end{example}

\begin{example}
	任意一个首一 $n$ 次多项式
	\[
	f(\lambda)=\lambda^n+a_1\lambda^{n-1}+\cdots+a_{n-1}\lambda+a_n
	\]
	都是某个矩阵的特征多项式.事实上，取友阵
	\[
	C_f=\begin{pmatrix}
	0&0&\cdots&0&-a_n\\
	1&0&\cdots&0&-a_{n-1}\\
	0&1&\cdots&0&-a_{n-2}\\
	\vdots&\vdots&\ddots&\vdots&\vdots\\
	0&0&\cdots&1&-a_1
\end{pmatrix},
\]
则
\[
f_{C_f}(\lambda)=f(\lambda).
\]
\end{example}

\begin{lemma}
设 $C_f$ 为上面的伴随矩阵.若 $B\in M_n(\F)$ 且 $BC_f=C_fB$，则存在 $g(x)\in\F[x]$，使得
\[
B=g(C_f).
\]
\end{lemma}

\begin{proof}
记标准基为 $e_1,\cdots,e_n$.由于
\[
e_i=C_f^{i-1}e_1,
\qquad 1\leq i\leq n,
\]
故 $e_1,C_f e_1,\cdots,C_f^{n-1} e_1$ 是$V$的一个基，设
\[
Be_1=b_0e_1+b_1C_fe_1+\cdots+b_{n-1}C_f^{n-1}e_1.
\]
令 $g(x)=b_0+b_1x+\cdots+b_{n-1}x^{n-1}$.因 $B$ 与 $C_f$ 可交换，故对每个 $i$，
\[
Be_i=B C_f^{i-1}e_1=C_f^{i-1}Be_1=g(C_f)e_i.
\]
于是 $B=g(C_f)$.
\end{proof}
设 $A\in M_n(\F)$.求 $A$ 的特征值与特征向量通常按如下步骤：
\begin{enumerate}
	\item 求特征多项式
	\[
	f_A(\lambda)=\det(\lambda E_n-A).
	\]
	由 $f_A(\lambda)=0$ 得到特征值 $\lambda_1,\cdots,\lambda_s$.
	\item 对每个特征值 $a$，解齐次线性方程组
	\[
	(A-aE_n)x=0.
	\]
	其解空间 $\mathcal N(A-aE_n)$ 中的非零向量就是属于 $a$ 的特征向量.
\end{enumerate}
需要注意：一般的初等行变换或初等列变换会改变特征值


设 $V$ 为 $n$ 维线性空间，$T\in\End(V)$，取 $V$ 的一组基
\[
\alpha_1,\cdots,\alpha_n,
\]
并设 $T$ 在这组基下的矩阵为 $A\in M_n(\F)$，即
\[
\bigl(T\alpha_1,\cdots,T\alpha_n\bigr)
=\bigl(\alpha_1,\cdots,\alpha_n\bigr)A.
\]
对 $\xi=(\alpha_1,\cdots,\alpha_n)x\in V$，其中 $x\in\F^n$，有
\[
T\xi
=\bigl(T\alpha_1,\cdots,T\alpha_n\bigr)x
=\bigl(\alpha_1,\cdots,\alpha_n\bigr)Ax.
\]
因此当 $x\neq0$ 时，
\[
T\xi=a\xi
\Longleftrightarrow
Ax=ax.
\]
于是 
\begin{corollary}
$a$ 是 $T$ 的特征值当且仅当 $a$ 是矩阵 $A$ 的特征值
\end{corollary}
并且坐标映射给出同构
\[
V_a(T)=\ker(T-a\id_V)
\cong
\mathcal N(A-aE_n).
\]
换言之我们可以通过考虑$T$在某个基下的矩阵$A$的特征值以考虑$T$的特征值
\begin{example}
设
\[
V=\langle 1,\sin x,\cos x,e^x\rangle,
\qquad D=\frac{d}{dx}.
\]
在基 $1,\sin x,\cos x,e^x$ 下，$D$ 的矩阵为
\[
[D]=
\begin{pmatrix}
	0&0&0&0\\
	0&0&-1&0\\
	0&1&0&0\\
	0&0&0&1
\end{pmatrix}.
\]
因此
\[
f_D(\lambda)=\lambda(\lambda^2+1)(\lambda-1).
\]
在 $\mathbb R$ 上，$D$ 的特征值为 $0,1$；在 $\mathbb C$ 上，还多出 $i,-i$.
\end{example}
根据代数基本定理，我们有
\begin{corollary}
	有限维复线性空间上的任意线性变换必然有复特征值与复特征向量
\end{corollary}
\begin{corollary}
设 $V$ 是有限维复线性空间，$S,T\in\End(V)$，且 $ST=TS$.则 $S,T$ 有公共的特征向量.
\end{corollary}

\begin{proof}
取 $S$ 的一个特征值 $a$.因为 $ST=TS$，所以 $S$ 的特征子空间
\[
V_a(S)=\ker(S-aE)
\]
是 $T$-子空间.由于 $V_a(S)\neq0$ 且是复线性空间，$T|_{V_a(S)}$ 有特征向量.该向量同时也是 $S$ 的特征向量.
\end{proof}

\begin{corollary}
设 $A,B\in M_n(\mathbb C)$ 且 $AB=BA$.则 $A,B$ 在 $\mathbb C^n$ 中有公共的特征向量.
\end{corollary}


下面把不变子空间与矩阵形状联系起来.

\begin{lemma}
设 $T\in\End(V)$，$\dim_{\F}V=n$.则 $T$ 有一个 $r$ 维不变子空间，当且仅当存在 $V$ 的一组基，使得 $T$ 在这组基下的矩阵形如
\[
\begin{pmatrix}
	A_{11}&A_{12}\\
	0&A_{22}
\end{pmatrix},
\qquad A_{11}\in M_r(\F).
\]
这种矩阵称为准上三角矩阵.
\end{lemma}

\begin{proof}
若 $U$ 是一个 $r$ 维 $T$-子空间，取 $U$ 的一组基
\[
\varepsilon_1,\cdots,\varepsilon_r
\]
并扩充为 $V$ 的一组基
\[
\varepsilon_1,\cdots,\varepsilon_r,
\varepsilon_{r+1},\cdots,\varepsilon_n.
\]
由于 $T(\varepsilon_i)\in U$，$1\leq i\leq r$，故矩阵左下角为 $0$.反过来，若矩阵具有上述形状，则前 $r$ 个基向量张成的子空间就是 $T$-子空间.
\end{proof}

\begin{corollary}
设 $T\in\End(V)$，$\dim_{\F}V=n$.则 $V$ 可分解为若干个 $T$-子空间的直和，当且仅当存在 $V$ 的一组基，使得 $T$ 在这组基下的矩阵为准对角阵
\[
\begin{pmatrix}
	A_1&&&0\\
	&A_2&&\\
	&&\ddots&\\
	0&&&A_s
\end{pmatrix}.
\]
\end{corollary}

\begin{proof}
若
\[
V=V_1\oplus\cdots\oplus V_s
\]
且每个 $V_i$ 都是 $T$-子空间，则分别取 $V_i$ 的一组基并合并为 $V$ 的一组基.此时 $T$ 在该基下的矩阵就是准对角阵.反过来，准对角阵的每个对角块对应一组基向量张成的 $T$-子空间.
\end{proof}
\begin{corollary}
	设$T \in \End(V),$且$\dim_{\F} V = n.$则 \(T\) 可对角化当且仅当 \(V\) 有一个基，使得 \(T\) 在这个基下的矩阵是对角阵.
	
	此时，这个基的每个向量都是 \(T\) 的特征向量.
	也就是说，$	T \text{ 可对角化}$当且仅当
	$T \text{ 有 } n \text{ 个线性无关的特征向量}.$
\end{corollary}
\section{相似}

有限维空间上的线性变换在不同基下的矩阵是相似的.设$\dim_{\F}V=n,
 T\in\End(V),$
取 $V$ 的基 $\alpha_1,\ldots,\alpha_n$与$\beta_1,\ldots,\beta_n$，并设
\[
(\alpha_1,\ldots,\alpha_n)=(\beta_1,\ldots,\beta_n)P,
\qquad P\in \GL_n(\F),
\]
设$T$在两基下的矩阵为$A,B$:
\[
(T\alpha_1\ldots T\alpha_n)
=(\alpha_1\ldots\alpha_n)A,
(T\beta_1\ldots T\beta_n)=(\beta_1\ldots\beta_n)B
\]
根据我们在第四章中的结果，有
\[
A=P^{-1}BP.
\]

\begin{definition}
	设 $A,B\in M_n(\F)$.若存在 $P\in \GL_n(\F)$，使得
	\[
	A=P^{-1}BP,
	\]
	则称 $A$ 与 $B$ 相似，记作 $A\sim B$.
\end{definition}

显然，相似关系是等价关系

\begin{proposition}
	设 $A,B\in M_n(\F)$.则 $A\sim B$ 当且仅当 $A,B$ 分别是同一个 $n$ 维线性空间上同一个线性变换在两组不同基下的矩阵.
\end{proposition}

\begin{example}
	若 $A\sim B$，则对任意 $k\geq 1$，有
	\[
	A^k\sim B^k.
	\]
	事实上，若 $A=P^{-1}BP$，则
	\[
	A^k=(P^{-1}BP)^k=P^{-1}B^kP.
	\]
	更一般地，对任意 $h(x)\in\F[x]$，有
	\[
	h(A)=P^{-1}h(B)P,
	\]
	所以 $h(A)\sim h(B)$.
\end{example}

\begin{proposition}
	设 $A,B\in M_n(\F)$ 且 $A\sim B$.则
	\begin{enumerate}
		\item $\rankA=\rankB$；
		\item 对任意 $c\in\F$，
		\[
		\dim\ker(A-cE_n)=\dim\ker(B-cE_n);
		\]
		\item $A,B$ 有相同的特征多项式，从而$\operatorname{tr}A = \operatorname{tr}B,|A|=|B|$：
		\item $A,B$ 有相同的极小多项式：
	\end{enumerate}
\end{proposition}

\begin{proof}
	设 $A=P^{-1}BP$.
	
	1. 由于左右乘可逆矩阵不改变秩，故 $\rankA=\rankB$.
	
	2. 有
	\[
	A-cE_n=P^{-1}(B-cE_n)P.
	\]
	映射 $x\mapsto Px$ 给出
	\[
	\ker(A-cE_n)\cong\ker(B-cE_n),
	\]
	从而维数相等.
	
	3. 由行列式乘法性，
	\[
	\begin{aligned}
		\det(\lambda E_n-A)
		&=\det\bigl(P^{-1}(\lambda E_n-B)P\bigr) \\
		&=\det(\lambda E_n-B).
	\end{aligned}
	\]
	4.设$A$的极小多项式为$m_A $,$B$的极小多项式为$m_B $,由$A\sim B$,得$m_B(A)\sim m_B(B)=0$从而$m_B $是$A$的零化多项式，$m_B|m_A$,同理$m_A|m_B$,故$m_A=m_B$
\end{proof}
\begin{example}
	$A,B\in M_n(\F)$则$A\sim B\rightarrow A'\sim B'$
\end{example}

\begin{example}
	两个 $n$ 阶对角矩阵相似当且仅当它们的对角元仅差一个排列.即
	\[
	D_1=\operatorname{diag}(a_1,\ldots,a_n),\qquad
	D_2=\operatorname{diag}(b_1,\ldots,b_n),
	\]
	则
	\[
	D_1\sim D_2
	\iff
	(b_1,\ldots,b_n) \text{ 与 } (a_1,\ldots,a_n) \text{ 差一个排列}.
	\]	例如
	\[
	\begin{pmatrix}
		1&&&0\\
		&1&&\\
		&&\ddots&\\
		0&&&3
	\end{pmatrix}.
	\text{与}
	\begin{pmatrix}
		1&&&0\\
		&3&&\\
		&&\ddots&\\
		0&&&1
	\end{pmatrix}.
	\]
\end{example}

\begin{proof}
	若 $D_1\sim D_2$，则二者有相同的特征多项式.于是
	\[
	f_{D_1}(\lambda)=\prod_{i=1}^n(\lambda-a_i),
	\qquad
	f_{D_2}(\lambda)=\prod_{i=1}^n(\lambda-b_i).
	\]
	由特征多项式相同可知，每个特征值出现的重数相同，故
	\[
	\{a_1,\ldots,a_n\}=\{b_1,\ldots,b_n\}
	\]
	作为含重数的集合相等.于是 $(b_1,\ldots,b_n)$ 只差 $(a_1,\ldots,a_n)$ 的一个排列.
	
	反过来，若存在排列 $\sigma\in S_n$，使得
	\[
	b_i=a_{\sigma(i)},\qquad 1\leq i\leq n,
	\]
	取置换矩阵
	\[
	P=(e_{\sigma(1)},e_{\sigma(2)},\ldots,e_{\sigma(n)}).
	\]
	则对任意 $1\leq i\leq n$，有
	\[
	D_1Pe_i=	D_1e_{\sigma(i)}
	=
	a_{\sigma(i)}e_{\sigma(i)}
	=
	b_iPe_i
	=
	PD_2e_i.
	\]
	因此
	\[
	D_1P=PD_2.
	\]
	由于 $P$ 可逆，所以
	\[
	P^{-1}D_1P=D_2.
	\]
	故 $D_1\sim D_2$.
\end{proof}
\begin{example}
	若 $A_1\sim B_1$，$A_2\sim B_2$，则
	\[
	\begin{pmatrix}
		A_1&0\\0&A_2
	\end{pmatrix}
	\sim
	\begin{pmatrix}
		B_1&0\\0&B_2
	\end{pmatrix}.
	\]
	事实上，若 $A_i=P_i^{-1}B_iP_i$，$i=1,2$，则取
	\[
	P=\begin{pmatrix}
		P_1&0\\0&P_2
	\end{pmatrix},
	\]
	即可.多个分块的情形同理.
\end{example}

\begin{proposition}
	设 $T\in\End(V)$，$\dim_{\F}V=n$，$W$ 是 $T$ 的一个不变子空间，$\dim W=r$.取 $W$ 的一组基
	\[
	\varepsilon_1,\ldots,\varepsilon_r,
	\]
	并扩充为 $V$ 的一组基
	\[
	\varepsilon_1,\ldots,\varepsilon_r,\varepsilon_{r+1},\ldots,\varepsilon_n.
	\]
	则 $T$ 在这组基下的矩阵形如
	\[
	\begin{pmatrix}
		A&C\\
		0&B
	\end{pmatrix},
	\]
	其中 $A$ 是 $T|_W$ 在基 $\varepsilon_1,
	\ldots,\varepsilon_r$ 下的矩阵.
\end{proposition}

\begin{proof}
	由于 $W$ 是 $T$-子空间，对 $1\leq i\leq r$ 有
	\[
	T\varepsilon_i\in W=\langle \varepsilon_1,\ldots,\varepsilon_r\rangle.
	\]
	故矩阵左下角分块为 $0$，即得到准上三角阵.
\end{proof}

\begin{corollary}
	在上述记号下，若 $T$ 的矩阵为
	\[
	\begin{pmatrix}
		A&C\\
		0&B
	\end{pmatrix},
	\]
	则
	\[
	f_T(\lambda)=f_A(\lambda)f_B(\lambda).
	\]
	特别地，若
	\[
	V=V_1\oplus\cdots\oplus V_s
	\]
	且每个 $V_i$ 都是 $T$-子空间，则取各 $V_i$ 的基合并为 $V$ 的基，$T$ 的矩阵为准对角阵
	\[
	\begin{pmatrix}
		A_1&&0\\
		&\ddots&\\
		0&&A_s
	\end{pmatrix},
	\]
	并且
	\[
	f_T(\lambda)=f_{A_1}(\lambda)
	\cdots f_{A_s}(\lambda)
	=f_{T|_{V_1}}(\lambda)
	\cdots f_{T|_{V_s}}(\lambda).
	\]
\end{corollary}

\begin{proposition}
	设 $A,B\in M_n(\F)$，若存在扩域 $\mathbb{K}\supseteq\F$，使得 $A,B$ 在 $\mathbb{K}$ 上相似，则 $A,B$ 在 $\F$ 上相似.
\end{proposition}

\begin{proof}
	在 $\mathbb{K}$ 上相似等价于存在 $Q\in \GL_n(\mathbb{K})$，使得
	\[
	AQ=QB.
	\]
	考虑齐次线性方程组
	\[
	AX=XB,
	\qquad X=(x_{ij}).
	\]
	其系数都在 $\F$ 中.$AX=XB$是一个包含了$n^2$个未知数的齐次线性方程组，在$\mathbb{K}$上有解，而有解不依赖于数域，故在$\F$上有解
	
	设
	\[
	\mathcal W_{\F}=\{X\in M_n(\F):AX=XB\},
	\]
	取其一组 $\F$-基 $X_1,
	\ldots,X_r$.由于方程系数在 $\F$ 中，同一组矩阵在 $\mathbb{K}$ 上张成
	\[
	\mathcal W_{\mathbb{K}}=\{X\in M_n(K):AX=XB\}.
	\]
	于是存在 $c_1,
	\ldots,c_r\in \mathbb{K}$，使得
	\[
	Q=c_1X_1+\cdots+c_rX_r.
	\]
	令
	\[
	g(t_1,
	\ldots,t_r)=\det(t_1X_1+
	\cdots+t_rX_r)\in\F[t_1,
	\ldots,t_r].
	\]
	由于 $g(c_1,
	\ldots,c_r)=\det Q\neq0$，故 $g$ 不是零多项式.所以存在 $a_1,
	\ldots,a_r\in\F$，使得
	\[
	g(a_1,
	\ldots,a_r)\neq0.
	\]
	于是
	\[
	P=a_1X_1+\cdots+a_rX_r\in \GL_n(\F),
	\qquad AP=PB.
	\]
	因此 $A=PBP^{-1}$，即 $A\sim B$.
\end{proof}



\section{三角化与可对角化}

本节继续讨论特征值、特征向量、三角化以及可对角化.以下默认 $A\in M_n(\F)$，若 $\mathbb K\supseteq\F$ 是扩域，则可把 $A$ 看成 $M_n(\mathbb K)$ 中的矩阵.

\begin{lemma}
	设 $a\in\mathbb K$ 是 $A$ 的特征值，记
	\[
	V_a=\ker(A-aE_n),\qquad r=\dim_{\mathbb K}V_a.
	\]
	则 $A$ 在 $\mathbb K$ 上相似于形如
	\[
	\begin{pmatrix}
		aE_r& *\\
		0&B
	\end{pmatrix}
	\]
	的准上三角矩阵.因此
	\[
	f_A(\lambda)=(\lambda-a)^r f_B(\lambda),
	\]
	从而特征值 $a$ 的几何重数不超过它的代数重数.
\end{lemma}

\begin{proof}
	取 $V_a$ 的一组基 $\xi_1,\ldots,\xi_r$，并扩充为 $\mathbb K^n$ 的一组基
	\[
	\xi_1,\ldots,\xi_r,\xi_{r+1},\ldots,\xi_n.
	\]
	由于 $A\xi_i=a\xi_i$，$1\leq i\leq r$，故 $A$ 在这组基下的矩阵左下角为 $0$，左上角为 $aE_r$.于是特征多项式由准上三角矩阵的对角块相乘得到.
\end{proof}

\begin{proposition}
	(上三角化)
	设 $A\in M_n(\F)$.则 $A$ 的特征多项式在 $\F$ 上完全分解(也即特征值全属于$\F$)当且仅当 $A$ 在 $\F$ 上相似于一个上三角矩阵.
\end{proposition}

\begin{proof}
	必要性：若
	\[
	A\sim
	\begin{pmatrix}
		b_{11}&*&\cdots&*\\
		0&b_{22}&\cdots&*\\
		\vdots&\vdots&\ddots&\vdots\\
		0&0&\cdots&b_{nn}
	\end{pmatrix},
	\]
	则
	\[
	f_A(\lambda)=(\lambda-b_{11})\cdots(\lambda-b_{nn})
	\]
	在 $\F$ 上完全分解.
	
	反过来，对 $n$ 归纳.$n=1$ 显然.设 $n>1$，由于 $f_A$ 在 $\F$ 上完全分解，$A$ 有特征值 $b\in\F$.取非零特征向量 $\xi$，并扩充为一组基，则
	\[
	A\sim
	\begin{pmatrix}
		b&*\\
		0&A_1
	\end{pmatrix},
	\qquad A_1\in M_{n-1}(\F).
	\]
	由
	\[
	f_A(\lambda)=(\lambda-b)f_{A_1}(\lambda)
	\]
	知 $f_{A_1}$ 也在 $\F$ 上完全分解.由归纳假设，$A_1$ 相似于上三角矩阵，因此 $A$ 相似于上三角矩阵.
\end{proof}

\begin{corollary}\label{相似于一个上三角矩阵.}
	任意 $A\in M_n(\mathbb C)$ 都在 $\mathbb C$ 上相似于一个上三角矩阵.
\end{corollary}

\begin{corollary}
	设 $A\in M_n(\F)$ 的全部特征值为 $\lambda_1,\ldots,\lambda_n$，按代数重数计.则对任意 $g(x)\in\F[x]$，$g(A)$ 的全部特征值为
	\[
	g(\lambda_1),\ldots,g(\lambda_n),
	\]
	并且
	\[
	\det g(A)=g(\lambda_1)\cdots g(\lambda_n),
	\qquad
	\operatorname{tr} g(A)=g(\lambda_1)+\cdots+g(\lambda_n).
	\]
\end{corollary}

\begin{proof}
	由上三角化，存在可逆矩阵 $P$ 使
	\[
	P^{-1}AP=U
	\]
	为上三角矩阵，且 $U$ 的对角元为 $\lambda_1,\ldots,\lambda_n$.于是 $g(U)$ 仍为上三角矩阵，其对角元为 $g(\lambda_1),\ldots,g(\lambda_n)$.又 $g(A)\sim g(U)$，结论成立.
\end{proof}

\begin{definition}
	设 $A\in M_n(\F)$，$\mathbb K\supseteq\F$.若存在 $P\in \GL_n(\mathbb K)$，使得
	\[
	P^{-1}AP
	\]
	为对角矩阵，则称 $A$ 在 $\mathbb K$ 上可对角化.若 $\mathbb K=\F$，则简称 $A$ 可对角化.
\end{definition}

显然，若 $A\sim B$，则 $A$ 可对角化当且仅当 $B$ 可对角化.并且若 $A$ 在 $\F$ 上可对角化，则在任意扩域 $\mathbb K\supseteq\F$ 上仍可对角化.


	设 $T\in\End(V)$，$\dim_{\F}V=n$,设$T$在某基下对应的矩阵是$A$.则下列条件等价：
	\begin{enumerate}
		\item $T$ 可对角化；
		\item $V$ 可分解为 $T$ 的特征子空间的直和；
		\item $V$ 有一组由 $T$ 的特征向量组成的基.
		\item $A$在$\F$上可对角化
	\end{enumerate}
	此时，若
	\[
	V=V_{a_1}\oplus\cdots\oplus V_{a_s},
	\qquad \dim V_{a_i}=t_i,
	\]
	并取 $V_{a_i}$ 的基
	\[
	\xi_{i1},\ldots,\xi_{it_i},
	\]
	则在基
	\[
	\xi_{11},\ldots,\xi_{1t_1},\ldots,
	\xi_{s1},\ldots,\xi_{st_s}
	\]
	下，$T$ 的矩阵为
	\[
	\begin{pmatrix}
		a_1E_{t_1}&&&0\\
		&a_2E_{t_2}&&\\
		&&\ddots&\\
		0&&&a_sE_{t_s}
	\end{pmatrix}.
	\]

\begin{proposition}
	设 $A\in M_n(\F)$，$\mathbb K\supseteq\F$.则 $A$ 在 $\mathbb K$ 上可对角化当且仅当 $\mathbb K^n$ 有一组由 $A$ 的特征向量组成的基.
\end{proposition}

\begin{proof}
	若 $P^{-1}AP=D=\operatorname{diag}(\lambda_1,\ldots,\lambda_n)$，则
	\[
	AP=PD.
	\]
	令 $P=(\xi_1,
	\ldots,\xi_n)$，则
	\[
	A\xi_i=\lambda_i\xi_i,
	\qquad 1\leq i\leq n.
	\]
	由于 $P$ 可逆，$\xi_1,\ldots,\xi_n$ 是 $\mathbb K^n$ 的一组基.反过来，若取到这样的特征向量基，以它们为列组成 $P$，则 $P^{-1}AP$ 为对角矩阵.
\end{proof}

\begin{example}
	矩阵
	\[
	A=\begin{pmatrix}0&1\\-1&0\end{pmatrix}
	\]
	在 $\mathbb R$ 上不可对角化，因为它没有实特征值；在 $\mathbb C$ 上可对角化，特征值为 $i,-i$.
\end{example}

\begin{example}
	设 $A=(a_{ij})\in M_n(\F)$ 是上三角矩阵，且
	\[
	a_{11}=\cdots=a_{nn}=a,
	\]
	若 $A\neq aE_n$，则 $A$ 在任意扩域上都不可对角化.
\end{example}

\begin{proof}
	$A$ 的唯一特征值为 $a$.若 $A$ 可对角化，则它相似于对角矩阵，而这个对角矩阵只能是 $aE_n$.但标量矩阵与任意矩阵可交换，且
	\[
	P^{-1}(aE_n)P=aE_n,
	\]
	故 $A=aE_n$，矛盾.
\end{proof}

\begin{theorem}
	设 $A\in M_n(\F)$，$\mathbb K\supseteq\F$，$m_A(x)$ 为 $A$ 的极小多项式.则 $A$ 在 $\mathbb K$ 上可对角化当且仅当 $m_A(x)$ 在 $\mathbb K$ 上可分解为互不相同的一次因式的乘积.
\end{theorem}

\begin{proof}
	若 $A$ 在 $\mathbb K$ 上可对角化，设
	\[
	A\sim D=\operatorname{diag}(\lambda_1E_{r_1},\ldots,\lambda_sE_{r_s}),
	\]
	其中 $\lambda_1,
	\ldots,\lambda_s$ 两两不同.则
	\[
	m_A(x)=m_D(x)=(x-\lambda_1)\cdots(x-\lambda_s).
	\]
	
	反过来，若
	\[
	m_A(x)=(x-a_1)\cdots(x-a_s),
	\qquad a_i\neq a_j\ (i\neq j),
	\]
	则
	\[
	(A-a_1E_n)\cdots(A-a_sE_n)=0.
	\]
	由准素分解定理，
	\[
	\mathbb K^n=\ker(A-a_1E_n)\oplus\cdots\oplus\ker(A-a_sE_n),
	\]
	故 $A$ 在 $\mathbb K$ 上可对角化.
\end{proof}
\begin{corollary}
	设 $A\in M_n(\F)$，$\mathbb K\supseteq\F$，$h(x)$ 为 $A$ 的零化多项式，若$(h,h')=1$,则$A$ 在 $\mathbb C$ 上可对角化
\end{corollary}
\begin{example}
	若 $A\in M_n(\mathbb R)$ 且
	\[
	A^4+E_n=0,
	\]
	则 $A$ 在 $\mathbb R$ 上不可对角化，但在 $\mathbb C$ 上可对角化.
\end{example}

\begin{proof}
	若 $A$ 在 $\mathbb R$ 上可对角化，则它的对角元应为实特征值.但 $x^4+1$ 在 $\mathbb R$ 上无根，故无实特征值，矛盾.
	
	在 $\mathbb C$ 上，$x^4+1$ 可分解为互不相同的一次因式，且 $m_A(x)\mid x^4+1$，故 $m_A$ 也没有重根，于是 $A$ 在 $\mathbb C$ 上可对角化.
\end{proof}

\begin{lemma}
	设 $A\in M_n(\F)$，$\mathbb K\supseteq\F$.则 $A$ 在 $\mathbb K$ 上可对角化当且仅当 $A$ 的特征值都属于 $\mathbb K$，且每个特征值的几何重数等于它的代数重数.
\end{lemma}

\begin{proof}
	若 $A$ 可对角化，则
	\[
	A\sim \operatorname{diag}(a_1E_{r_1},\ldots,a_sE_{r_s}),
	\qquad a_i\neq a_j.
	\]
	于是
	\[
	f_A(\lambda)=(\lambda-a_1)^{r_1}\cdots(\lambda-a_s)^{r_s},
	\]
	且属于 $a_i$ 的几何重数为 $r_i$，同时 $a_i$ 的代数重数也为 $r_i$
	
	反过来，若
	\[
	f_A(\lambda)=(\lambda-a_1)^{r_1}\cdots(\lambda-a_s)^{r_s},
	\qquad a_i\neq a_j,
	\]
	并且
	\[
	\dim\ker(A-a_iE_n)=r_i,
	\qquad 1\leq i\leq s,
	\]
	则不同特征值对应的特征子空间和为直和，且维数之和为
	\[
	r_1+\cdots+r_s=n.
	\]
	故
	\[
	\mathbb K^n=V_{a_1}\oplus\cdots\oplus V_{a_s},
	\]
	从而 $A$ 在 $\mathbb K$ 上可对角化.
\end{proof}

\begin{algo}
	判断 $A\in M_n(\F)$ 在扩域 $\mathbb K$ 上是否可对角化，通常按如下步骤：
	\begin{enumerate}
		\item 求
		\[
		f_A(\lambda)=\det(\lambda E_n-A).
		\]
		若 $f_A$ 在 $\mathbb K$ 上不能完全分解，则 $A$ 在 $\mathbb K$ 上不可对角化.
		\item 设 $a_1,
		\ldots,a_s$ 是 $f_A$ 在 $\mathbb K$ 中互不相同的根.对每个 $a_i$ 解
		\[
		(A-a_iE_n)x=0,
		\]
		其解空间维数为
		\[
		r_i=n-\rank(a_iE_n-A).
		\]
		若某个 $r_i$ 小于 $a_i$ 的代数重数，则不可对角化.
		\item 若每个特征值的几何重数都等于代数重数，取
		\[
		\ker(A-a_iE_n)
		\]
		的一组基 $\xi_{i1},\ldots,\xi_{ir_i}$，令
		\[
		P=(\xi_{11},\ldots,\xi_{1r_1},\xi_{21},\ldots,\xi_{sr_s}).
		\]
		则
		\[
		P^{-1}AP=
		\begin{pmatrix}
			a_1E_{r_1}&&&0\\
			&a_2E_{r_2}&&\\
			&&\ddots&\\
			0&&&a_sE_{r_s}
		\end{pmatrix}.
		\]
	\end{enumerate}
\end{algo}


\section{Hamilton-Cayley 定理}


\begin{example}
	先看可对角化情形.设 $A\in M_n(\F)$ 可对角化，且
\[
\F^n=V_{a_1}\oplus\cdots\oplus V_{a_s},
\qquad r_i=\dim V_{a_i}.
\]
于是 $A$ 的特征多项式为
\[
f_A(\lambda)=(\lambda-a_1)^{r_1}\cdots(\lambda-a_s)^{r_s}.
\]
对任意 $\alpha\in\F^n$，可唯一写成
\[
\alpha=\alpha_1+\cdots+\alpha_s,
\qquad \alpha_i\in V_{a_i}.
\]
因此
\[
\begin{aligned}
	f_A(A)\alpha
	&=(A-a_1E_n)^{r_1}\cdots(A-a_sE_n)^{r_s}
	(\alpha_1+\cdots+\alpha_s)       \\
	&=\sum_{i=1}^s
	\prod_{j=1}^s (A-a_jE_n)^{r_j}\alpha_i=0,
\end{aligned}
\]
因为在第 $i$ 项中含有因子 $(A-a_iE_n)^{r_i}\alpha_i=0$.故 $f_A(A)=0$.

也可以从矩阵形式看.若
\[
A\sim D=\begin{pmatrix}
	a_1E_{r_1}&&&0\\
	&a_2E_{r_2}&&\\
	&&\ddots&\\
	0&&&a_sE_{r_s}
\end{pmatrix},
\]
则
\[
f_A(D)=\prod_{i=1}^s (D-a_iE_n)^{r_i}=0,
\]
从而 $f_A(A)=0$.
\end{example}

\begin{example}
设 $A\in M_n(\F)$ 是幂零矩阵.若
\[
A\sim U,
\]
其中 $U\in M_n(\mathbb C)$ 是对角元全为 $0$ 的上三角矩阵，则 $U^n=0$，于是
\[
A^n\sim U^n=0.
\]
此时 $A$ 的特征多项式为
\[
f_A(\lambda)=\lambda^n,
\]
故 $f_A(A)=A^n=0$.
\end{example}

\begin{theorem}(Hamilton-Cayley 定理)
设 $A\in M_n(\F)$，其特征多项式为
\[
f_A(\lambda)=\det(\lambda E_n-A),
\]
则
\[
f_A(A)=0.
\]
换言之，每个方阵的特征多项式是他的零化多项式.
\end{theorem}

\begin{proof}
由\ref{相似于一个上三角矩阵.},$A$ 在$\C$相似于上三角矩阵.故只需证明上三角矩阵情形.

对 $n$ 归纳.$n=1$ 显然.设
\[
B=\begin{pmatrix}
b&*\\
0&B_1
\end{pmatrix},
\qquad B_1\in M_{n-1}(\F).
\]
则
\[
f_B(\lambda)=(\lambda-b)f_{B_1}(\lambda).
\]
由归纳假设，$f_{B_1}(B_1)=0$，所以
\[
f_{B_1}(B)=
\begin{pmatrix}
f_{B_1}(b)&*\\
0&0
\end{pmatrix}.
\]
于是
\[
\begin{aligned}
f_B(B)
&=(B-bE_n)f_{B_1}(B)\\
&=
\begin{pmatrix}
	0&*\\
	0&B_1-bE_{n-1}
\end{pmatrix}
\begin{pmatrix}
	f_{B_1}(b)&*\\
	0&0
\end{pmatrix}=0.
\end{aligned}
\]
因此上三角情形成立，从而一般情形也成立.
\end{proof}

\begin{proposition}
设 $T\in\End(V)$，且 $\dim V<\infty$.则 $T$ 的特征多项式是 $T$ 的一个零化多项式.
\end{proposition}

\begin{proof}
取 $V$ 的一组基，使 $T$ 在该基下的矩阵为 $A$.由 Hamilton-Cayley 定理，$f_A(A)=0$，故 $f_T(T)=0$.
\end{proof}

\begin{corollary}
设 $A\in M_n(\F)$，$m_A(x)$ 为 $A$ 的极小多项式，$f_A(x)$ 为 $A$ 的特征多项式，则
\[
m_A(x)\mid f_A(x).
\]
特别地，$A$ 的特征值一定是 $m_A(x)$ 的根.
\end{corollary}

\begin{example}
若
\[
f_A(\lambda)=(\lambda^2+1)(\lambda-1)^2,
\]
则极小多项式可能为
\[
(\lambda^2+1)(\lambda-1)
\qquad\text{或}\qquad
(\lambda^2+1)(\lambda-1)^2.
\]

\end{example}

\begin{example}
	
设 $A\in M_n(\F)$，则
\[
\F[A]=\{h(A):h(x)\in\F[x]\}
\]
满足
\[
\dim \F[A]\leq n.
\]
事实上，$\dim \F[A]$ 等于 $A$ 的极小多项式的次数.

\end{example}

\begin{proof}
由 Hamilton-Cayley 定理，$f_A(A)=0$.对任意 $h(x)\in\F[x]$，用带余除法写成
\[
h(x)=q(x)f_A(x)+r(x),
\qquad \deg r<n.
\]
于是
\[
h(A)=r(A)\in\langle E_n,A,\ldots,A^{n-1}\rangle,
\]
故 $\dim\F[A]\leq n$.
更精确地，$\dim\F[A]$ 等于 $A$ 的极小多项式的次数.
\end{proof}




\begin{example}
设 $A\in M_n(\F)$，$f_A(\lambda)=\det(\lambda E_n-A)$，$g(x)\in\F[x]$.则
\[
g(A)\in \GL_n(\F)
\iff
(f_A,g)=1.
\]
\end{example}

\begin{proof}
若 $(f_A,g)=1$，则存在 $u(x),v(x)\in\F[x]$，使
\[
u(x)f_A(x)+v(x)g(x)=1.
\]
代入 $A$ 得
\[
v(A)g(A)=E_n,
\]
故 $g(A)$ 可逆.

反之，若 $(f_A,g)\neq1$，设公共根为 $\lambda_0$.因为 $\lambda_0$ 是 $f_A$ 的根，故存在非零向量 $\xi$ 使
\[
A\xi=\lambda_0\xi.
\]
又 $g(\lambda_0)=0$，于是
\[
g(A)\xi=g(\lambda_0)\xi=0.
\]
故 $g(A)$ 不可逆，矛盾
\end{proof}

\begin{example}
设 $A,B\in M_n(\F)$，若 $A,B$ 的特征多项式互素，则矩阵方程
\[
AX=XB
\]
只有零解.
\end{example}

\begin{proof}
若 $AX=XB$，则对任意 $h(x)\in\F[x]$，有
\[
h(A)X=Xh(B).
\]
取 $h=f_A$，由 Hamilton-Cayley 定理得
\[
0=f_A(A)X=Xf_A(B).
\]
由于 $(f_A,f_B)=1$，由上例知 $f_A(B)$ 可逆，故 $X=0$.
\end{proof}

\begin{example}
设
\[
A=\begin{pmatrix}
A_1&&&0\\
&A_2&&\\
&&\ddots&\\
0&&&A_s
\end{pmatrix}
\]
是准对角矩阵，且 $A_1,\ldots,A_s$ 的特征多项式两两互素.若矩阵 $B$ 与 $A$ 可交换，则 $B$ 也是与 $A$ 具有相同分块方式的准对角矩阵.
\end{example}

\begin{proof}
按照 $A$ 的分块方式写
\[
B=(B_{ij})_{1\leq i,j\leq s}.
\]
由 $AB=BA$ 得
\[
A_iB_{ij}=B_{ij}A_j,
\qquad 1\leq i,j\leq s.
\]
当 $i\neq j$ 时，$A_i,A_j$ 的特征多项式互素，由上一例知 $B_{ij}=0$.故 $B$ 为相同分块方式的准对角矩阵.
\end{proof}

下面由 Hamilton-Cayley 定理，把准素分解定理应用于特征多项式.

\begin{example}
设 $A\in\End(V)$，$\dim_{\F}V<\infty$，$f(\lambda)$ 为 $A$ 的特征多项式.若
\[
f(\lambda)=g(\lambda)h(\lambda),
\qquad (g,h)=1,
\]
其中 $g(\lambda),h(\lambda)\in\F[\lambda]$ 均首一，则
\[
V=\ker g(A)\oplus\ker h(A).
\]
设
\[
A_1=A|_{\ker g(A)},
\qquad
A_2=A|_{\ker h(A)},
\]
则 $A_1$ 的特征多项式为 $g(\lambda)$，$A_2$ 的特征多项式为 $h(\lambda)$.
\end{example}

\begin{proof}
由 Hamilton-Cayley 定理，$f(A)=0$.又 $(g,h)=1$，由准素分解定理，
\[
V=\ker g(A)\oplus\ker h(A).
\]
记 $A_1=A|_{\ker g(A)}$，$A_2=A|_{\ker h(A)}$.由于两个核都是 $A$-不变子空间，故在适当基下
\[
A\sim \begin{pmatrix}
A_1&0\\
0&A_2
\end{pmatrix}.
\]
从而
\[
f_A(\lambda)=f_{A_1}(\lambda)f_{A_2}(\lambda).
\]
又 $g(A_1)=0$，故 $f_{A_1}$ 的根只能来自 $g$；同理，$f_{A_2}$ 的根只能来自 $h$.结合
\[
f_A(\lambda)=g(\lambda)h(\lambda),
\]
即得
\[
f_{A_1}(\lambda)=g(\lambda),
\qquad
f_{A_2}(\lambda)=h(\lambda).
\]
\end{proof}
\section{$\lambda$-矩阵}

我们已经知道，若 $A,B\in M_n(\F)$ 且 $A\sim B$，则
\[
\det(\lambda E_n-A)=\det(\lambda E_n-B).
\]
但仅有特征多项式相同一般不能推出 $A\sim B$.为了研究相似关系，需要引入 $\lambda$-矩阵.

\begin{definition}
	设 $A\in M_n(\F)$，称
	\[
	\lambda E_n-A
	\]
	为 $A$ 的特征矩阵.
	一般地，若
	\[
	A(\lambda)=\bigl(a_{ij}(\lambda)\bigr)\in M_n(\F[\lambda]),
	\]
	则称 $A(\lambda)$ 为一个 $n$ 阶 $\lambda$-矩阵.
\end{definition}

设 $A(\lambda)=\bigl(a_{ij}(\lambda)\bigr)\in M_n(\F[\lambda])$.记
\[
\deg A(\lambda)=\max_{1\leq i,j\leq n}\deg a_{ij}(\lambda),
\]
并约定 $\deg0=-\infty$.
若 $\deg A(\lambda)=m$，则可写成
\[
A(\lambda)=\lambda^mA_m+\lambda^{m-1}A_{m-1}+\cdots+\lambda A_1+A_0,
\qquad A_i\in M_n(\F).
\]

$M_n(\F[\lambda])$ 在通常的矩阵加法与乘法下成为一个环.若
\[
A(\lambda)=\bigl(a_{ij}(\lambda)\bigr),
\qquad
B(\lambda)=\bigl(b_{ij}(\lambda)\bigr),
\]
则
\[
A(\lambda)+B(\lambda)=\bigl(a_{ij}(\lambda)+b_{ij}(\lambda)\bigr),
\]
并且若 $C(\lambda)=A(\lambda)B(\lambda)$，则
\[
C(\lambda)=\bigl(c_{ij}(\lambda)\bigr),
\qquad
c_{ij}(\lambda)=\sum_{k=1}^n a_{ik}(\lambda)b_{kj}(\lambda).
\]
于是有
\[
\deg(A(\lambda)+B(\lambda))\leq \max\{\deg A(\lambda),\deg B(\lambda)\},
\]
\[
\deg(A(\lambda)B(\lambda))\leq \deg A(\lambda)+\deg B(\lambda).
\]


\begin{lemma}
	设
	\[
	A(\lambda)=\sum_{i=0}^m\lambda^iA_i,
	\qquad
	B(\lambda)=\sum_{j=0}^s\lambda^jB_j.
	\]
	则
	\[
	A(\lambda)B(\lambda)
	=
	\sum_{t=0}^{m+s}\lambda^t
	\left(\sum_{i+j=t}A_iB_j\right).
	\]
\end{lemma}

\begin{proof}
	直接由矩阵乘法与多项式乘法展开即可.
\end{proof}

\begin{lemma}
	设 $B\in M_n(\F)$.对任意 $P(\lambda)\in M_n(\F[\lambda])$，存在 $H(\lambda)\in M_n(\F[\lambda])$ 与 $R_1\in M_n(\F)$，使得
	\[
	P(\lambda)=(\lambda E_n-B)H(\lambda)+R_1.
	\]
	对任意 $Q(\lambda)\in M_n(\F[\lambda])$，存在 $G(\lambda)\in M_n(\F[\lambda])$ 与 $R_2\in M_n(\F)$，使得
	\[
	Q(\lambda)=G(\lambda)(\lambda E_n-B)+R_2.
	\]
\end{lemma}

\begin{proof}
	只证第一个式子.设
	\[
	P(\lambda)=\lambda^mP_m+\lambda^{m-1}P_{m-1}+\cdots+\lambda P_1+P_0.
	\]
	令
	\[
	H(\lambda)=\lambda^{m-1}H_{m-1}+\lambda^{m-2}H_{m-2}+\cdots+H_0.
	\]
	由
	\[
	P(\lambda)=(\lambda E_n-B)H(\lambda)+R_1
	\]
	比较各次项系数，得到
	\[
	\left\{
	\begin{aligned}
		H_{m-1}&=P_m,\\
		H_{m-2}&=P_{m-1}+BH_{m-1},\\
		&\ \ \vdots\\
		H_0&=P_1+BH_1,\\
		R_1&=P_0+BH_0.
	\end{aligned}
	\right.
	\]
	故可依次确定 $H_{m-1},\ldots,H_0,R_1$.
	第二个式子同理，只需从右乘以 $\lambda E_n-B$.
\end{proof}


\begin{definition}
	设 $A(\lambda)=\bigl(a_{ij}(\lambda)\bigr)\in M_n(\F[\lambda])$.定义
	\[
	\det A(\lambda)
	=
	\sum_{i_1\cdots i_n\in S_n}(-1)^{\tau(i_1\cdots i_n)}
	a_{1i_1}(\lambda)a_{2i_2}(\lambda)\cdots a_{ni_n}(\lambda)
	\in\F[\lambda].
	\]
\end{definition}

例如，若 $A\in M_n(\F)$，则
\[
\det(\lambda E_n-A)
\]
就是 $A$ 的特征多项式，是一个首一的 $n$ 次多项式.

$\lambda$-矩阵的行列式与普通矩阵的行列式具有相同性质.特别地：
\[
\det(A(\lambda)B(\lambda))=\det A(\lambda)\det B(\lambda),
\]
并且若 $A(\lambda)^*$ 表示 $A(\lambda)$ 的伴随矩阵，则
\[
A(\lambda)A(\lambda)^*=A(\lambda)^*A(\lambda)=\det A(\lambda)E_n.
\]

\begin{proposition}(用$\lambda$阵证明Hamilton-Cayley 定理)
	设 $A\in M_n(\F)$，
	\[
	f_A(\lambda)=\det(\lambda E_n-A).
	\]
	则
	\[
	f_A(A)=0.
	\]
\end{proposition}

\begin{proof}
	设
	\[
	f_A(\lambda)=\lambda^n+a_1\lambda^{n-1}+\cdots+a_{n-1}\lambda+a_n.
	\]
	由伴随矩阵公式，
	\[
	(\lambda E_n-A)^*(\lambda E_n-A)=f_A(\lambda)E_n.
	\]
	记
	\[
	(\lambda E_n-A)^*
	=
	\lambda^{n-1}B_{n-1}+\lambda^{n-2}B_{n-2}+\cdots+\lambda B_1+B_0,
	\qquad B_i\in M_n(\F).
	\]
	比较上式两边各次幂系数，得
	\[
	\left\{
	\begin{aligned}
		B_{n-1}&=E_n,\\
		B_{n-2}-B_{n-1}A&=a_1E_n,\\
		B_{n-3}-B_{n-2}A&=a_2E_n,\\
		&\ \ \vdots\\
		B_0-B_1A&=a_{n-1}E_n,\\
		-B_0A&=a_nE_n.
	\end{aligned}
	\right.
	\]
	于是
	\[
	\left\{
	\begin{aligned}
		B_{n-1}A^n&=A^n,\\
		B_{n-2}A^{n-1}-B_{n-1}A^n&=a_1A^{n-1},\\
		B_{n-3}A^{n-2}-B_{n-2}A^{n-1}&=a_2A^{n-2},\\
		&\ \ \vdots\\
		B_0A-B_1A^2&=a_{n-1}A,\\
		-B_0A&=a_nE_n.
	\end{aligned}
	\right.
	\]
	将这些等式相加，左边望远相消，得到
	\[
	A^n+a_1A^{n-1}+\cdots+a_{n-1}A+a_nE_n=0.
	\]
	即 $f_A(A)=0$.
\end{proof}

\begin{definition}
	设 $A(\lambda)\in M_n(\F[\lambda])$.若 $A(\lambda)$ 的非零子式的最高阶数为 $r$，则称 $A(\lambda)$ 的秩为 $r$，记为
	\[
	r(A(\lambda))=r
	\qquad\text{或}\qquad
	\rankA(\lambda)=r.
	\]
\end{definition}

显然
\[
0\leq r(A(\lambda))\leq n.
\]

\begin{corollary}
	设 $A(\lambda)\in M_n(\F[\lambda])$.则
	\[
	\det A(\lambda)\neq0
	\iff
	r(A(\lambda))=n.
	\]
\end{corollary}

\begin{example}
	设 $A\in M_n(\F)$，则
	\[
	r(\lambda E_n-A)=n.
	\]
	因为 $\det(\lambda E_n-A)$ 是首一 $n$ 次多项式，故不为零.
\end{example}
\section{$\lambda$矩阵的可逆性}

\begin{definition}[可逆 $\lambda$-矩阵]
	设 $A(\lambda)\in M_n(\F[\lambda])$.若存在 $B(\lambda)\in M_n(\F[\lambda])$，使得
	\[
	B(\lambda)A(\lambda)=A(\lambda)B(\lambda)=E_n,
	\]
	则称 $A(\lambda)$ 是可逆的，并称 $B(\lambda)$ 为 $A(\lambda)$ 的逆，记为
	\[
	A(\lambda)^{-1}=B(\lambda).
	\]
\end{definition}

	逆矩阵若存在则唯一.事实上，若 $B_1(\lambda),B_2(\lambda)$ 都是 $A(\lambda)$ 的逆，则
	\[
	B_1(\lambda)=B_1(\lambda)E_n
	=B_1(\lambda)A(\lambda)B_2(\lambda)
	=E_nB_2(\lambda)=B_2(\lambda).
	\]

\begin{lemma}
	设 $A(\lambda)\in M_n(\F[\lambda])$.则 $A(\lambda)$ 可逆当且仅当
	\[
	\det A(\lambda)\in \F\setminus\{0\}.
	\]
\end{lemma}

\begin{proof}
	若 $A(\lambda)$ 可逆，则存在 $B(\lambda)$ 使得
	\[
	A(\lambda)B(\lambda)=E_n.
	\]
	于是
	\[
	\det A(\lambda)\det B(\lambda)=1.
	\]
	故 $\det A(\lambda)$ 是 $\F[\lambda]$ 中的可逆元，所以
	\[
	\det A(\lambda)\in\F\setminus\{0\}.
	\]
	反过来，若 $\det A(\lambda)=c\in\F\setminus\{0\}$，则由Laplace定理
	\[
	A(\lambda)A(\lambda)^*=A(\lambda)^*A(\lambda)=cE_n
	\]
	可知
	\[
	A(\lambda)^{-1}=c^{-1}A(\lambda)^*\in M_n(\F[\lambda]).
	\]
\end{proof}

\begin{corollary}
	设 $A(\lambda),B(\lambda)\in M_n(\F[\lambda])$.则 $A(\lambda)B(\lambda)$ 可逆当且仅当 $A(\lambda)$ 与 $B(\lambda)$ 都可逆.
\end{corollary}

注意：对 $A\in M_n(\F)$，特征矩阵 $\lambda E_n-A$ 一般不可逆，因为
\[
\det(\lambda E_n-A)
\]
通常不是 $\F[\lambda]$ 中的可逆元.但它的秩为 $n$，即
\[
r(\lambda E_n-A)=n.
\]
\begin{corollary}
	$A(\lambda)\in M_n(\F[\lambda])$,则$A(\lambda)$可逆$\iff$$\exists B(\lambda)\in M_n(\F[\lambda])$使$A(\lambda)B(\lambda)=E_n $或$B(\lambda)A(\lambda)=E_n $
\end{corollary}

\begin{definition}[相抵]
	设 $A(\lambda),B(\lambda)\in M_n(\F[\lambda])$.若存在可逆 $\lambda$-矩阵 $P(\lambda),Q(\lambda)$，使得
	\[
	P(\lambda)A(\lambda)Q(\lambda)=B(\lambda),
	\]
	则称 $A(\lambda)$ 与 $B(\lambda)$ 相抵,
	这是一个等价关系.
\end{definition}

\begin{proposition}
	设 $A,B\in M_n(\F)$.则
	\[
	A\text{ 与 }B\text{ 相似}
	\iff
	\lambda E_n-A\text{ 与 }\lambda E_n-B\text{ 相抵}.
	\]
\end{proposition}

\begin{proof}
	若 $A\sim B$，则存在 $C\in \GL_n(\F)$，使得
	\[
	B=CAC^{-1}.
	\]
	于是
	\[
	C(\lambda E_n-A)C^{-1}=\lambda E_n-B,
	\]
	故 $\lambda E_n-A$ 与 $\lambda E_n-B$ 相抵.
	
	反过来，若存在可逆 $P(\lambda),Q(\lambda)\in M_n(\F[\lambda])$，使得
	\[
	P(\lambda)(\lambda E_n-A)Q(\lambda)=\lambda E_n-B.
	\]
	由 $\lambda$阵的带余除法，存在 $R_1,R_2\in M_n(\F)$ 以及 $H(\lambda),G(\lambda)\in M_n(\F[\lambda])$，使得
	\[
	P(\lambda)=(\lambda E_n-B)H(\lambda)+R_1,
	\qquad
	Q(\lambda)=G(\lambda)(\lambda E_n-B)+R_2.
	\]
	代入上式并整理，可得
	\[
	R_1(\lambda E_n-A)R_2
	=(\lambda E_n-B)\bigl(E_n-L(\lambda)(\lambda E_n-B)\bigr)
	\]
	其中 $L(\lambda)\in M_n(\F[\lambda])$.比较两边关于 $\lambda$ 的次数，得 $L(\lambda)=0$.故
	\[
	R_1(\lambda E_n-A)R_2=\lambda E_n-B.
	\]
	比较 $\lambda$ 的系数与常数项，得
	\[
	R_1R_2=E_n,
	\qquad
	R_1AR_2=B.
	\]
	因此 $R_1\in \GL_n(\F)$，且 $B=R_1AR_1^{-1}$，所以 $A\sim B$.
\end{proof}

\begin{definition}[初等 $\lambda$-变换]
	$\F$ 上的 $\lambda$-矩阵的初等行变换为下列三种操作：
	\begin{enumerate}
		\item 交换两行，记为$(i,j)$；
		\item 用 $\F$ 中的非零数乘某一行,记为$c(i)$；
		\item 将 $\F[\lambda]$ 中的一个多项式倍的某一行加到另一行上,记为$(j) + \varphi (\lambda)(i)$.
	\end{enumerate}
	初等列变换同理定义.相应的初等矩阵称为初等 $\lambda$-矩阵，分别记为$P(i,j),P(c(i)),P((j)+\varphi(
	\lambda)(i))$
	
	初等列变换与初等行变换统称为$\lambda$-矩阵的初等变换
\end{definition}


	初等行变换等价于左乘初等 $\lambda$-矩阵，初等列变换等价于右乘初等 $\lambda$-矩阵.初等 $\lambda$-矩阵均可逆，且逆矩阵仍是初等 $\lambda$-矩阵.

	容易发现，初等$\lambda$-矩阵可逆，其逆仍为初等$\lambda$-矩阵
	\[P(i,j)P(i,j) =E_n,P(c(i))P(c^{-1}(i)) = E_n ,P((j)+\varphi(
	\lambda)(i))P((j)-\varphi(
	\lambda)(i)) = E_n\]
\begin{theorem}(Smith 标准形)
	设 $0\neq A(\lambda)\in M_n(\F[\lambda])$.则 $A(\lambda)$ 可经过有限次初等行变换与初等列变换化为
	\[
	\begin{pmatrix}
		d_1(\lambda) & & & &&& \\
		& d_2(\lambda) & & & &&\\
		& & \ddots & & &&\\
		& & & d_r(\lambda) &&& \\
		& & & & 0&&\\
		&&&&&\ddots&\\
		&&&&&&0\\
	\end{pmatrix},
	\]
	其中 $d_i(\lambda)$ 均为首一多项式，并且
	\[
	d_1(\lambda)\mid d_2(\lambda)\mid\cdots\mid d_r(\lambda).
	\]
	上式称为 $A(\lambda)$ 的 Smith 标准形.
\end{theorem}

\begin{proof}
	对 $n$ 归纳。$n=1$ 时显然成立。设 $n>1$，并设结论对 $n-1$ 阶矩阵成立。
	
	记
	\[
	\mathcal X=\{Y(\lambda)\in M_n(\F[\lambda]):Y(\lambda)
	\text{ 可由 }A(\lambda)\text{ 经有限步初等变换得到}\}.
	\]
	若 $A(\lambda)=0$，结论显然。否则取
	\[
	Z(\lambda)=(z_{ij}(\lambda))\in\mathcal X
	\]
	使得 $z_{11}(\lambda)\neq0$，$z_{11}(\lambda)$ 为首一，且
	\[
	\deg z_{11}(\lambda)
	\]
	在所有可作为 $(1,1)$ 元的非零首一多项式中最小。
	
	先证
	\[
	z_{11}(\lambda)\mid z_{i1}(\lambda),\qquad
	z_{11}(\lambda)\mid z_{1j}(\lambda),
	\qquad 1\leq i,j\leq n.
	\]
	若 $z_{11}\nmid z_{i1}$，由带余除法，
	\[
	z_{i1}=z_{11}q+r,\qquad r\neq0,\quad \deg r<\deg z_{11}.
	\]
	将第 $i$ 行减去第 $1$ 行的 $q$ 倍，再交换第 $1$ 行与第 $i$ 行，
	便得到一个仍属于 $\mathcal X$ 的矩阵，其 $(1,1)$ 元为 $r$。
	这与 $z_{11}$ 的选取矛盾。故 $z_{11}\mid z_{i1}$。
	同理，对列作初等变换可得 $z_{11}\mid z_{1j}$。
	
	于是由初等行、列变换可把第一列和第一行除 $(1,1)$ 元外全部化为 $0$，即
	\[
	Z(\lambda)\sim
	\begin{pmatrix}
		z_{11}(\lambda) & 0 & \cdots & 0\\
		0 & z'_{22}(\lambda) & \cdots & z'_{2n}(\lambda)\\
		\vdots & \vdots & \ddots & \vdots\\
		0 & z'_{n2}(\lambda) & \cdots & z'_{nn}(\lambda)
	\end{pmatrix}.
	\]
	对右下角的 $(n-1)$ 阶矩阵用归纳假设，可得
	\[
	A(\lambda)\sim
	\operatorname{diag}\bigl(z_{11}(\lambda),d_2(\lambda),\ldots,d_n(\lambda)\big),
	\]
	其中
	\[
	d_2(\lambda)\mid d_3(\lambda)\mid\cdots\mid d_n(\lambda).
	\]
	
	只需再证 $z_{11}(\lambda)\mid d_2(\lambda)$。若不然，由带余除法，
	\[
	d_2=z_{11}q_1+r_1,\qquad r_1\neq0,\quad \deg r_1<\deg z_{11}.
	\]
	对矩阵
	\[
	D=\operatorname{diag}\bigl(z_{11},d_2,d_3,\ldots,d_n\bigr)
	\]
	作如下初等变换：
	\[
	\begin{pmatrix}
		z_{11} & 0 & \cdots & 0\\
		0 & d_2 & \cdots & 0\\
		\vdots & \vdots & \ddots & \vdots\\
		0 & 0 & \cdots & d_n
	\end{pmatrix}
	\longrightarrow
	\begin{pmatrix}
		z_{11} & d_2 & \cdots & 0\\
		0 & d_2 & \cdots & 0\\
		\vdots & \vdots & \ddots & \vdots\\
		0 & 0 & \cdots & d_n
	\end{pmatrix}
	\longrightarrow
	\begin{pmatrix}
		z_{11} & r_1 & \cdots & 0\\
		0 & d_2 & \cdots & 0\\
		\vdots & \vdots & \ddots & \vdots\\
		0 & 0 & \cdots & d_n
	\end{pmatrix}.
	\]
	再交换前两列，得到
	\[
	\begin{pmatrix}
		r_1 & z_{11} & \cdots & 0\\
		d_2 & 0 & \cdots & 0\\
		\vdots & \vdots & \ddots & \vdots\\
		0 & 0 & \cdots & d_n
	\end{pmatrix}\in\mathcal X.
	\]
	其 $(1,1)$ 元为 $r_1$，且 $\deg r_1<\deg z_{11}$，与 $z_{11}$ 的最小性矛盾。
	所以
	\[
	z_{11}\mid d_2.
	\]
	因此
	\[
	z_{11}\mid d_2\mid d_3\mid\cdots\mid d_n,
	\]
	从而得到所需的 Smith 标准型。若某些 $d_i=0$，把非零项排在前面即可。
\end{proof}
\section{$\lambda$-矩阵的行列式因子、不变因子与初等因子}


\begin{definition}
	
设 $X(\lambda)\in M_n(\F[\lambda])$，且 $r(X)=r\geq 1$. 对 $1\leq k\leq r$，
称 $X(\lambda)$ 所有 $k$ 阶子式的首一最大公因式为 $X(\lambda)$ 的第 $k$ 个行列式因子，记为
\[
D_k(\lambda).
\]
若 $k>r$，则记
\[
D_k(\lambda)=0.
\]

\end{definition}
由Laplace定理 我们可以得到以下的推论、


\begin{corollary}
		设 $X(\lambda)\in M_n(\F[\lambda])$ 的全部非零行列式因子为
	\[
	D_1(\lambda),D_2(\lambda),\ldots,D_r(\lambda),
	\]
	则
	\[
	D_i(\lambda)\mid D_{i+1}(\lambda),\qquad 1\leq i<r.
	\]
	特别地，若 $D_k(\lambda)=1$，则
	\[
	D_1(\lambda)=\cdots=D_{k-1}(\lambda)=1.
	\]
\end{corollary}


\begin{proof}
	由 Laplace 展开，每个 $(i+1)$ 阶子式都可以写成若干个 $i$ 阶子式的 $\F[\lambda]$-线性组合.
	故 $D_i(\lambda)$ 整除每个 $(i+1)$ 阶子式，从而
	\[
	D_i(\lambda)\mid D_{i+1}(\lambda).
	\]
	后一句由整除链立即得到.
\end{proof}

\begin{example}
	设 $0\neq X(\lambda)=(x_{ij}(\lambda))\in M_n(\F[\lambda])$，则 $X$ 的第一个行列式因子就是所有 $x_{ij}(\lambda)$ 的首一最大公因式.
\end{example}

\begin{example}
	设 $A\in M_n(\F)$，则 $\lambda E_n-A$ 的第 $n$ 个行列式因子为
	\[
	D_n(\lambda)=\det(\lambda E_n-A),
	\]
	即 $A$ 的特征多项式. 此时
	\[
	r(\lambda E_n-A)=n.
	\]
\end{example}

\begin{example}
	考虑友矩阵
	\[
	A=\begin{pmatrix}
		0&0&0&\cdots&0&-a_n\\
		1&0&0&\cdots&0&-a_{n-1}\\
		0&1&0&\cdots&0&-a_{n-2}\\
		\vdots&\vdots&\vdots&\ddots&\vdots&\vdots\\
		0&0&0&\cdots&0&-a_2\\
		0&0&0&\cdots&1&-a_1
	\end{pmatrix}.
	\]
	则 $\lambda E_n-A$ 的全部行列式因子为
	\[
	D_1=\cdots=D_{n-1}=1,
	\]
	且
	\[
	D_n=\lambda^n+a_1\lambda^{n-1}+\cdots+a_{n-1}\lambda+a_n.
	\]
\end{example}

\begin{example}
	设 $A\in M_n(\F)\subset M_n(\F[\lambda])$，且 $r(A)=r>0$. 则 $A$ 的全部非零行列式因子为
	\[
	D_1=\cdots=D_r=1.
	\]
	这是因为 $A$ 有一个非零的 $r$ 阶子式，而该子式是 $\F$ 中的非零数.
\end{example}

\begin{example}
	若 $X(\lambda)\in M_n(\F[\lambda])$ 有一个 $r$ 阶子式是非零常数 $c\in\F$，则
	\[
	D_1=\cdots=D_r=1.
	\]
\end{example}

\begin{example}
	设 $X(\lambda)\in M_n(\F[\lambda])$，则 $X(\lambda)$ 可逆当且仅当它的全部行列式因子为
	\[
	D_1=\cdots=D_n=1.
	\]
\end{example}

\begin{example}
	若
	\[
	X(\lambda)\sim
	\begin{pmatrix}
		d_1(\lambda) \\
		& d_2(\lambda) \\
		&& \ddots \\
		&&& d_r(\lambda) \\
		&&&& 0 \\
		&&&&& \ddots
	\end{pmatrix},
	\]
	其中 $d_i(\lambda)$ 都是首一的，且
	\[
	d_i(\lambda)\mid d_{i+1}(\lambda),\qquad 1\leq i<r,
	\]
	则 $X(\lambda)$ 的行列式因子为
	\[
	D_1=d_1,
	\]
	\[
	D_2=d_1d_2,
	\]
	\[
	\cdots
	\]
	\[
	D_r=d_1d_2\cdots d_r,
	\qquad
	D_j=0,\quad r+1\leq j\leq n.
	\]
	因此
	\[
	d_1=D_1,
	\qquad
	d_i=\frac{D_i}{D_{i-1}},\quad 2\leq i\leq r.
	\]
\end{example}

\begin{lemma}
	$\lambda$-矩阵的初等变换不改变它的行列式因子.
\end{lemma}

\begin{proof}
	只需验证一次初等变换. 设 $X=(x_{ij})$ 经过一次初等行变换得到 $Y=(y_{ij})$.
	
	若是行交换，则每个子式至多差一个符号；若是某一行乘以非零常数，则每个子式至多差一个非零常数.
	
	若是将第 $i$ 行加上第 $j$ 行的 $\varphi(\lambda)$ 倍，则 $Y$ 的任意 $k$ 阶子式要么是 $X$ 的一个 $k$ 阶子式，
	要么是 $X$ 的两个 $k$ 阶子式的 $\F[\lambda]$-线性组合. 因此 $X$ 的第 $k$ 个行列式因子整除 $Y$ 的所有 $k$ 阶子式.
	反向作初等变换同理，故二者的第 $k$ 个行列式因子相同.
	
	列变换的证明完全相同.
\end{proof}


由 Smith 标准型和上面的例子可知，若 $0\neq X(\lambda)\in M_n(\F[\lambda])$ 的 Smith 标准型为
\[
\operatorname{diag}(d_1(\lambda),\ldots,d_r(\lambda),0,\ldots,0),
\]
则其中的 $d_i(\lambda)$ 被行列式因子唯一确定：
\[
d_1=D_1,
\qquad
d_i=\frac{D_i}{D_{i-1}},\quad 2\leq i\leq r.
\]

\begin{definition}
	设 $0\neq X(\lambda)\in M_n(\F[\lambda])$. 称其 Smith 标准型中的
	\[
	d_1(\lambda),\ldots,d_r(\lambda)
	\]
	为 $X(\lambda)$ 的不变因子.
	
	设 $A\in M_n(\F)$. 称 $\lambda E_n-A$ 的不变因子为 $A$ 的不变因子；称 $\lambda E_n-A$ 的行列式因子为 $A$ 的行列式因子.
\end{definition}

\begin{proposition}
	设 $X(\lambda),Y(\lambda)\in M_n(\F[\lambda])$. 则下列条件等价：
	\begin{enumerate}
		\item $X(\lambda)$ 与 $Y(\lambda)$ 相抵；
		\item $X(\lambda)$ 与 $Y(\lambda)$ 的行列式因子相同；
		\item $X(\lambda)$ 与 $Y(\lambda)$ 的不变因子相同；
		\item $X(\lambda)$ 与 $Y(\lambda)$ 有相同的 Smith 标准型.
	\end{enumerate}
\end{proposition}

\begin{proof}
	相抵矩阵可以由初等变换互相得到，故行列式因子相同.
	行列式因子唯一确定不变因子，不变因子唯一确定 Smith 标准型.
	反过来，若有相同 Smith 标准型，则二者都与同一个 Smith 标准型相抵，故彼此相抵.
\end{proof}
立刻得到以下定理

\begin{theorem}
	设 $A,B\in M_n(\F)$，则
	\[
	A\sim B
	\]
	当且仅当 $A$ 与 $B$ 的不变因子相同，当且仅当 $A$ 与 $B$ 的行列式因子相同.
\end{theorem}

\begin{corollary}
	设 $A,B\in M_n(\F)$，$K$ 是数域且 $\F\subset K$. 则
	\[
	A\text{ 在 }\F\text{ 上相似于 }B
	\quad\Longleftrightarrow\quad
	A\text{ 在 }K\text{ 上相似于 }B.
	\]
\end{corollary}
因为行列式因子不依赖于数域
\begin{corollary}
	设 $A\in M_n(\F)$，则
	\[
	A\sim A'.
	\]
\end{corollary}

\begin{proof}
	因为
	\[
	\lambda E_n-A'=(\lambda E_n-A)',
	\]
	而转置不改变行列式因子，所以 $A$ 与 $A'$ 的行列式因子相同.
\end{proof}
\begin{lemma}[同一不可约因子的重排]
	设 $p(\lambda)\in\F[\lambda]$ 为首一不可约多项式，且 $p\nmid g_1,\ p\nmid g_2$.
	则对任意 $a,b\geq 0$，有
	\[
	\operatorname{diag}\bigl(\cdots,p^a g_1,\cdots,\ p^b g_2,\cdots\bigr)
	\text{与}
	\operatorname{diag}\bigl(\cdots,p^b g_1,\cdots,\ p^a g_2,\cdots\bigr)\text{相抵}
	\]
\end{lemma}

\begin{proof}
	对行列式因子我们只需要考虑其最大公因子,不妨设$a>b$
	\[(p^ag_1,p^b g_2) = p^b(p^{a-b}g_1,p g_2)= p^b(g_1, g_2)\]
	\[(p^bg_1,p^a g_2) = p^b(g_1,p^{a-b} g_2)= p^b(g_1, g_2)\]
	从而$D_1,\cdots,D_n$相同
\end{proof}

由此可知，在对角 $\lambda$-矩阵中，对于同一个首一不可约因子 $p(\lambda)$，
可以把它在各对角元中的幂次重新排列，而不改变相抵关系.




\begin{definition}
	设 $X(\lambda)$ 的非零不变因子为
	\[
	d_1(\lambda),\ldots,d_r(\lambda).
	\]
	也即$X(\lambda)$的Smith标准型
	\[
	X(\lambda)\sim
	\begin{pmatrix}
		1\\
		&\ddots\\
		&&1\\
		&&&d_1(\lambda) \\
		&&&& d_2(\lambda) \\
		&&&&& \ddots \\
		&&&&&& d_r(\lambda) \\
		&&&&&&& 0 \\
		&&&&&&&& \ddots
	\end{pmatrix},
	\]
	把每个非单位不变因子在 $\F[\lambda]$ 中分解为首一不可约多项式幂的乘积：
	\[
	d_i(\lambda)=p_{i1}(\lambda)^{e_{i1}}\cdots p_{it_i}(\lambda)^{e_{it_i}}.
	\]
	其中出现的所有首一素因子幂
	\[
	p_{ij}(\lambda)^{e_{ij}}
	\]
	称为 $X(\lambda)$ 的初等因子. 对 $A\in M_n(\F)$，称 $\lambda E_n-A$ 的初等因子为 $A$ 的初等因子.
\end{definition}
\begin{example}
	若
	\[
	X=\operatorname{diag}\bigl(1,\lambda^2(\lambda-1),1,\lambda,(\lambda-1)^3\bigr),
	\]
	则 $X$ 的初等因子为
	\[
	\lambda,\quad \lambda^2,
	\qquad
	\lambda-1,
	\qquad
	(\lambda-1)^3.
	\]
	按同一素因子从小到大排列，再从右向左相乘，得到 Smith 标准型
	\[
	\operatorname{diag}\bigl(1,1,1,\lambda(\lambda-1),\lambda^2(\lambda-1)^3\bigr).
	\]
\end{example}


\begin{lemma}[满秩 $\lambda$-矩阵由初等因子确定相抵类]
	设 $X,Y\in M_n(\F[\lambda])$ 都是满秩的 $\lambda$-矩阵. 则
	\[
	X\sim Y
	\]
	当且仅当 $X,Y$ 的初等因子相同.
\end{lemma}

\begin{proof}
	若 $X\sim Y$，则 $X,Y$ 的不变因子相同，因此初等因子相同.
	
	反过来，只需说明满秩 $\lambda$-矩阵的不变因子可由全部初等因子唯一确定.
	设满秩矩阵 $Z\in M_n(\F[\lambda])$ 的全部初等因子按首一不可约因子分类为
	\[
	\begin{array}{cccc}
		p_1(\lambda)^{\ell_{11}}&p_1(\lambda)^{\ell_{12}}&\cdots&p_1(\lambda)^{\ell_{1t_1}}\\
		p_2(\lambda)^{\ell_{21}}&p_2(\lambda)^{\ell_{22}}&\cdots&p_2(\lambda)^{\ell_{2t_2}}\\
		\vdots&\vdots&&\vdots\\
		p_s(\lambda)^{\ell_{s1}}&p_s(\lambda)^{\ell_{s2}}&\cdots&p_s(\lambda)^{\ell_{st_s}}
	\end{array}
	\]
	其中 $p_1(\lambda),\ldots,p_s(\lambda)$ 在 $\F[\lambda]$ 上首一、不可约且两两不同，并且
	\[
	\ell_{i1}\geq \ell_{i2}\geq\cdots\geq \ell_{it_i}\geq1.
	\]
	若 $t_i<n$，就在第 $i$ 行后面补上 $n-t_i$ 个 $1$，得到一个 $s\times n$ 的表格：
	\[
	\begin{array}{ccccccc}
		p_1^{\ell_{11}}&p_1^{\ell_{12}}&\cdots&p_1^{\ell_{1t_1}}&1&\cdots&1\\
		p_2^{\ell_{21}}&p_2^{\ell_{22}}&\cdots&p_2^{\ell_{2t_2}}&1&\cdots&1\\
		\vdots&\vdots&&\vdots&\vdots&&\vdots\\
		p_s^{\ell_{s1}}&p_s^{\ell_{s2}}&\cdots&p_s^{\ell_{st_s}}&1&\cdots&1
	\end{array}
	\]
	其中不足处均视为指数 $0$.
	令每一列的乘积分别为
	\[
	d_n(\lambda),d_{n-1}(\lambda),\ldots,d_1(\lambda).
	\]
	由于同一行中指数递减，故
	\[
	d_1(\lambda)\mid d_2(\lambda)\mid\cdots\mid d_n(\lambda).
	\]
	于是
	\[
	\operatorname{diag}(d_1(\lambda),\ldots,d_n(\lambda))
	\]
	就是 $Z$ 的 Smith 标准型. 因而初等因子唯一确定不变因子. 由此得证.
\end{proof}
\begin{corollary}
	设 $X,Y\in M_n(\F[\lambda])$. 则下列条件等价：
	\begin{enumerate}
		\item $X, Y$相抵；
		\item $X$ 与 $Y$ 的不变因子相同；
		\item $X$ 与 $Y$ 的行列式因子相同；
		\item $X$ 与 $Y$ 的初等因子相同.
	\end{enumerate}
\end{corollary}
\begin{corollary}
		设 $A,B\in M_n(\F)$. 则下列条件等价：
	\begin{enumerate}
		\item $A\sim B$；
		\item $A$ 与 $B$ 的不变因子相同；
		\item $A$ 与 $B$ 的行列式因子相同；
		\item $A$ 与 $B$ 的初等因子相同.
	\end{enumerate}
\end{corollary}



	设 $X\in M_n(\F[\lambda])$ 是满秩的 $\lambda$-矩阵，且 $X$ 与对角 $\lambda$-矩阵
	\[
	Y=\operatorname{diag}(h_1(\lambda),h_2(\lambda),\ldots,h_n(\lambda))
	\]
	相抵，其中每个 $h_j(\lambda)$ 都首一. 则出现在所有 $h_j(\lambda)$ 的标准分解式中的首一素幂因子，按重数计算，正好是 $X$ 的全部初等因子.



	因 $X\sim Y$，$X$ 与 $Y$ 有相同的初等因子，故只需看 $Y$.
	设
	\[
	h_j(\lambda)=p_1(\lambda)^{e_{1j}}p_2(\lambda)^{e_{2j}}\cdots p_m(\lambda)^{e_{mj}},
	\qquad e_{ij}\geq0.
	\]
	按照上面引理的操作方式，就可以得到\[d_n(\lambda),\cdots, d_1(\lambda)\]
	那么，就可以得到$X$的Smith标准型
	
	也就是说，我们只要把$X$通过初等变换变为对角阵之后就可以使用上述的做法得到Smith标准型。
\begin{example}
	求下列 $X$ 的相抵标准型：
	\[
	X=\begin{pmatrix}
	1&&&&\\
	&\lambda^2(\lambda-1)&&&\\
	&&1&&\\
	&&&\lambda&\\
	&&&& (\lambda-1)^3
\end{pmatrix}.
\]
$X$ 的初等因子为
\[
\lambda,\quad \lambda^2,
\qquad
\lambda-1,
\qquad
(\lambda-1)^3.
\]
把同一首一素因子按次数从大到小排成列，再由右向左合并，得到
\[
\begin{array}{ccccc}
	\lambda^2&\lambda&1&1&1\\
	(\lambda-1)^3&\lambda-1&1&1&1
\end{array}
\]
故 $X$ 的相抵标准型为
\[
\begin{pmatrix}
	1&&&&\\
	&1&&&\\
	&&1&&\\
	&&&\lambda(\lambda-1)&\\
	&&&&\lambda^2(\lambda-1)^3
\end{pmatrix}.
\]
\end{example}

\begin{corollary}
设 $A\in M_n(\F)$，$\mathbb{K}$ 是数域且 $\F\subset \mathbb{K}$. 则 $A$ 在 $\mathbb{K}$ 上可对角化，当且仅当 $A$ 在 $\mathbb{K}$ 上的初等因子全为 $\mathbb{K}$ 上的一次多项式.
\end{corollary}


\begin{proof}
若 $A$ 在 $\mathbb{K}$ 上可对角化，则
\[
A\sim \operatorname{diag}(a_1,\ldots,a_n),\qquad a_i\in \mathbb{K}.
\]
于是 $A$ 的初等因子为
\[
\lambda-a_1,\ldots,\lambda-a_n.
\]
反过来，若 $A$ 的初等因子全为一次多项式
\[
\lambda-b_1,
\ldots,
\lambda-b_n,
\]
则 $\lambda E_n-A$ 与
\[
\operatorname{diag}(\lambda-b_1,
\ldots,
\lambda-b_n)
\]
相抵，故
\[
A\sim \operatorname{diag}(b_1,
\ldots,b_n).
\]
所以 $A$ 在 $\mathbb{K}$ 上可对角化.
\end{proof}

\begin{corollary}
对于任意准对角矩阵，把每个对角块的初等因子合起来，相同的要重复计，即可得到该准对角矩阵的全部初等因子.
\end{corollary}

\begin{proof}
设
\[
A=\begin{pmatrix}
	A_1&&&\\
	&A_2&&\\
	&&\ddots&\\
	&&&A_s
\end{pmatrix},
\qquad
A_i\in M_{r_i}(\F),
\quad r_1+\cdots+r_s=n.
\]
则
\[
\lambda E_n-A=
\begin{pmatrix}
	\lambda E_{r_1}-A_1&&&\\
	&\lambda E_{r_2}-A_2&&\\
	&&\ddots&\\
	&&&\lambda E_{r_s}-A_s
\end{pmatrix}.
\]
对每个对角块分别作 $\lambda$-矩阵的初等变换，设
\[
\lambda E_{r_i}-A_i\sim D_i(\lambda),
\qquad 1\leq i\leq s,
\]
其中 $D_i(\lambda)$ 为该块的 Smith 标准型. 于是
\[
\lambda E_n-A
\sim
\begin{pmatrix}
	D_1(\lambda)&&&\\
	&D_2(\lambda)&&\\
	&&\ddots&\\
	&&&D_s(\lambda)
\end{pmatrix}.
\]
故整体的初等因子就是所有块的初等因子合并所得，且相同的要重复计.
\end{proof}

\section{有理标准型与 Jordan 标准型}

设
\[
h(\lambda)=\lambda^m+b_1\lambda^{m-1}+\cdots+b_{m-1}\lambda+b_m\in\F[\lambda]
\]
为首一多项式.定义与 $h(\lambda)$ 对应的友矩阵
\[
\mathcal{R}(h)=
\begin{pmatrix}
	0&0&\cdots&0&-b_m\\
	1&0&\cdots&0&-b_{m-1}\\
	0&1&\cdots&0&-b_{m-2}\\
	\vdots&\vdots&\ddots&\vdots&\vdots\\
	0&0&\cdots&1&-b_1
\end{pmatrix}\in M_m(\F).
\]


\begin{lemma}
	设 $\mathcal{R}=\mathcal{R}(h)$，则
	\[
	\det(\lambda E_m-\mathcal{R})=h(\lambda),
	\]
	并且 $\lambda E_m-\mathcal{R}$ 的行列式因子为
	\[
	D_1=\cdots=D_{m-1}=1,\qquad D_m=h(\lambda).
	\]
	因此 $\mathcal{R}(h)$ 的唯一非平凡不变因子为 $h(\lambda)$.
\end{lemma}

\begin{proof}
	由伴随矩阵的定义直接计算可得
	\[
	\det(\lambda E_m-\mathcal{R})=h(\lambda).
	\]
	又在 $\lambda E_m-\mathcal{R}$ 中含有一个 $(m-1)$ 阶子式等于 $1$，故
	\[
	D_1=\cdots=D_{m-1}=1.
	\]
	最后
	\[
	D_m=\det(\lambda E_m-\mathcal{R})=h(\lambda).
	\]
\end{proof}

$\mathcal{R}$在$\F$上这只有一个初等因子当且仅当$h(\lambda)$是素幂，也即$h(\lambda)=p(\lambda)^\ell$，其中 $p(\lambda)$ 是首一不可约多项式

$\forall p(\lambda)^l,\exists \mathcal{R}\in M_{n}(\F)$使$\mathcal{R}$的初等因子恰为$p(\lambda)^\ell$,对相应的$\mathcal{R}$我们定义
\begin{definition}
	 设$p(\lambda)$ 是首一不可约多项式，且\[p(\lambda)^\ell = \lambda^m + b_1 \lambda^{m-1} + \cdots + b_m\]
	 ，令\[
	 \mathcal{R}(h)=
	 \begin{pmatrix}
	 	0&0&\cdots&0&-b_m\\
	 	1&0&\cdots&0&-b_{m-1}\\
	 	0&1&\cdots&0&-b_{m-2}\\
	 	\vdots&\vdots&\ddots&\vdots&\vdots\\
	 	0&0&\cdots&1&-b_1
	 \end{pmatrix}\in M_m(\F).
	 \]
	 称$\mathcal{R}$为$p(\lambda)^\ell$的有理块
\end{definition}
$\F$上的一个有理块不可能相似与一个由更低阶有理块构成的$\F$上的准对角阵

$\R$上的有理块只能对应$(\lambda -a)^k$与$(\lambda^2 + p\lambda +q)^k$;$\C$上的有理块只能对应$(\lambda -c)^k$
\begin{theorem}(有理标准型)
	设 $A\in M_n(\F)$，且 $A$ 在 $\F$ 上的全部初等因子为
	\[
	p_i(\lambda)^{\ell_{ij}},
	\qquad
	1\leq i\leq s,\quad 1\leq j\leq t_i,
	\]
	其中 $p_1(\lambda),\ldots,p_s(\lambda)$ 是 $\F[\lambda]$ 中两两不同的首一不可约多项式.
	令
	\[
	\mathcal{R}_{ij}=\mathcal{R}\bigl(p_i(\lambda)^{\ell_{ij}}\bigr).
	\]
	则 $A$ 相似于准对角矩阵
	\[
	\mathcal{R}=
	\operatorname{diag}
	\bigl(
	\mathcal{R}_{11},\ldots,\mathcal{R}_{1t_1},
	\mathcal{R}_{21},\ldots,\mathcal{R}_{2t_2},
	\ldots,
	\mathcal{R}_{s1},\ldots,\mathcal{R}_{st_s}
	\bigr).
	\]
	称 $\mathcal{R}$ 为 $A$ 在 $\F$ 上的一个有理标准型.
\end{theorem}

\begin{proof}
	对每个初等因子 $p_i(\lambda)^{\ell_{ij}}$，取对应有理块
	\[
	\mathcal{R}_{ij}=\mathcal{R}\bigl(p_i(\lambda)^{\ell_{ij}}\bigr).
	\]
	由上一个引理，$\mathcal{R}_{ij}$ 的唯一非平凡初等因子就是
	\[
	p_i(\lambda)^{\ell_{ij}}.
	\]
	故准对角矩阵 $\mathcal{R}$ 的全部初等因子正好是 $A$ 的全部初等因子.从而，$A\sim \mathcal{R}$.
\end{proof}

\begin{corollary}
	设 $A\in M_n(\F)$，$A$ 在 $\F$ 上的全部初等因子为
	\[
	p_i(\lambda)^{\ell_{ij}},
	\qquad
	1\leq i\leq s,\quad 1\leq j\leq t_i,
	\]
	其中同一个 $i$ 下可令
	\[
	\ell_{i1}\geq\ell_{i2}\geq\cdots\geq\ell_{it_i}.
	\]
	则 $A$ 的极小多项式为
	\[
	m_A(\lambda)
	=
	p_1(\lambda)^{\ell_{11}}
	p_2(\lambda)^{\ell_{21}}
	\cdots
	p_s(\lambda)^{\ell_{s1}},
	\]
	即每个不可约因子只取出现的最高次数.特别地，$A$ 的极小多项式等于 $A$ 的最后一个不变因子.
\end{corollary}

\begin{proof}
	有理标准型中，对应 $p_i$ 的各块的零化多项式为
	\[
	p_i(\lambda)^{\ell_{i1}},\ldots,p_i(\lambda)^{\ell_{it_i}}.
	\]
	\textbf{整个准对角矩阵的极小多项式为各块极小多项式的最小公倍式}，因此只取每个 $p_i$ 的最高幂.
\end{proof}
也就是说：
	$A$ 的特征多项式为
	\[
	f_A(\lambda)
	=
	\prod_{i=1}^{s}\prod_{j=1}^{t_i}
	p_i(\lambda)^{\ell_{ij}}.
	\]

\begin{algo}
	求 $A$ 在 $\F$ 上的有理标准型，可以按如下步骤：
	\begin{enumerate}
		\item 求出 $\lambda E_n-A$ 的全部初等因子；
		\item 对每个初等因子 $p_i(\lambda)^{\ell_{ij}}$ 写出对应的友矩阵 $\mathcal{R}_{ij}$；
		\item 将这些 $\mathcal{R}_{ij}$ 拼成准对角矩阵，即得 $A$ 的有理标准型.
	\end{enumerate}
\end{algo}

	

	设 $f(\lambda)\in\F[\lambda]$ 首一，且
	\[
	f(\lambda)=p(\lambda)^a
	\]
	为某个首一不可约多项式的幂.若
	\[
	a=a_1+\cdots+a_s,\qquad
	a_1\geq a_2\geq\cdots\geq a_s\geq1,
	\]
	则称 $(a_1,\ldots,a_s)$ 是 $a$ 的一个划分.相应地，
	\[
	p(\lambda)^{a_1},\ldots,p(\lambda)^{a_s}
	\]
	给出以 $f(\lambda)$ 为特征多项式的初等因子的一种可能.
	记 $a$ 的划分个数为 $\mathcal{P}(a)$.
\begin{proposition}
		设 $f(\lambda)\in\F[\lambda]$ 首一，且
	\[
	\deg f(\lambda)=n\geq1.
	\]
	设 $f(\lambda)$ 在 $\F$ 上的标准分解式为
	\[
	f(\lambda)=p_1(\lambda)^{r_1}\cdots p_s(\lambda)^{r_s},
	\qquad r_i\geq1.
	\]
	则 $\F$ 上以 $f(\lambda)$ 为特征多项式的所有 $n$ 阶方阵，在相似意义下可以分为
	\[
	\mathcal{P}(r_1)\mathcal{P}(r_2)\cdots \mathcal{P}(r_s)
	\]
	类.
\end{proposition}

\begin{example}
	设 $A\in M_n(\F)$，且 $A$ 的初等因子可分为两组：
	第一组形如 $p(\lambda)^k$，其中 $p(0)\neq0$；第二组形如 $\lambda^r$.
	则存在可逆矩阵 $P$，使
	\[
	P^{-1}AP=
	\begin{pmatrix}
		B&0\\
		0&C
	\end{pmatrix},
	\]
	其中 $B$ 可逆，$C$ 幂零.
\end{example}


\begin{definition}
	设 $a\in\F$，$r\in\mathbb{Z}^+$.定义 $r$ 阶 Jordan 块为
	\[
	J_r(a)=
	\begin{pmatrix}
		a&0&0&\cdots&0\\
		1&a&0&\cdots&0\\
		0&1&a&\cdots&0\\
		\vdots&\vdots&\ddots&\ddots&\vdots\\
		0&0&\cdots&1&a
	\end{pmatrix}.
	\]
	以Jordan块为对角块的准对角阵称为一个Jordan阵
\end{definition}
$J_r(a)$的行列式因子为
\[
D_1=\cdots=D_{r-1}=1,\qquad
D_r=(\lambda-a)^r.
\]
它与 $(\lambda-a)^r$ 的有理块相似
因此 $J_r(a)$ 的唯一非平凡初等因子为$(\lambda-a)^r.$


\begin{lemma}
	设
	\[
	J=\operatorname{diag}
	\bigl(
	J_{r_1}(a_1),\ldots,J_{r_s}(a_s)
	\bigr).
	\]
	则 $J$ 在 $\F$ 上的全部初等因子为
	\[
	(\lambda-a_1)^{r_1},\ldots,(\lambda-a_s)^{r_s},
	\]
	相同的要重复计算.
\end{lemma}


\begin{definition}
	设 $A\in M_n(\F)$。若 $\F$ 上的与 $A$ 相似的 Jordan 矩阵存在，
	则称它为 $A$ 在 $\F$ 上的 Jordan 标准型。
\end{definition}

\begin{theorem}[Jordan 标准型]
	设 $A\in M_n(\F)$。则 $A$ 与 $\F$ 上的一个 Jordan 矩阵相似，
	当且仅当 $A$ 的特征值全属于 $\F$。
	此时，与 $A$ 相似的 Jordan 矩阵在允许交换对角块顺序的意义下是唯一的。
\end{theorem}

\begin{proof}
	必要性,若
	\[
	A\sim J=\operatorname{diag}\bigl(J_{r_1}(a_1),\ldots,J_{r_s}(a_s)\bigr),
	\]
	则 $A$ 与 $J$ 的特征多项式相同，而
	\[
	f_J(\lambda)=(\lambda-a_1)^{r_1}\cdots(\lambda-a_s)^{r_s}.
	\]
	因此 $A$ 的特征值全属于 $\F$。
	
	反过来，设 $A$ 的特征值全属于 $\F$，即
	\[
	f_A(\lambda)=(\lambda-a_1)^{m_1}\cdots(\lambda-a_s)^{m_s},
	\]
	其中 $a_i\in\F$ 且两两不同。于是 $A$ 的全部初等因子必形如
	\[
	\begin{array}{cccc}
		(\lambda-a_1)^{r_{11}},&
		(\lambda-a_1)^{r_{12}},&
		\cdots,&
		(\lambda-a_1)^{r_{1t_1}},\\
		(\lambda-a_2)^{r_{21}},&
		(\lambda-a_2)^{r_{22}},&
		\cdots,&
		(\lambda-a_2)^{r_{2t_2}},\\
		\vdots&\vdots&&\vdots\\
		(\lambda-a_s)^{r_{s1}},&
		(\lambda-a_s)^{r_{s2}},&
		\cdots,&
		(\lambda-a_s)^{r_{st_s}} .
	\end{array}
	\]
	其中，对每个 $i$，$(r_{i1},r_{i2},\ldots,r_{it_i})$ 是 $m_i$ 的一个划分。
	对每个初等因子 $(\lambda-a_i)^{r_{ij}}$ 取 Jordan 块
	\[
	J_{r_{ij}}(a_i),\qquad 1\leq i\leq s,\quad 1\leq j\leq t_i.
	\]
	令
	\[
	J=\operatorname{diag}\bigl(
	J_{r_{11}}(a_1),\ldots,J_{r_{1t_1}}(a_1),
	\ldots,
	J_{r_{s1}}(a_s),\ldots,J_{r_{st_s}}(a_s)
	\bigr).
	\]
	由上一个引理，$J$ 与 $A$ 的初等因子完全相同，故 $A\sim J$。
	
	唯一性由初等因子的唯一性得到：每个初等因子
	$(\lambda-a_i)^{r_{ij}}$ 正好对应一个 Jordan 块
	$J_{r_{ij}}(a_i)$，因此只可能交换对角块的顺序。
\end{proof}

\begin{corollary}
	设 $A\in M_n(\F)$,在 $\mathbb C$ 上有 Jordan 标准型。
\end{corollary}
\begin{proposition}[由零度数恢复 Jordan 块]
	设 $A\in M_n(\F)$ 在 $\F$ 上有 Jordan 标准型，$\lambda\in\F$ 是 $A$ 的一个特征值.记
	\[
	\nu_k=\dim\ker(A-\lambda E)^k,\qquad k=0,1,2,\ldots,
	\]
	其中约定 $\nu_0=0$.
	
	若 $A$ 关于特征值 $\lambda$ 的 Jordan 块大小分别为
	\[
	s_1,s_2,\ldots,s_m,
	\]
	即 $\lambda$ 对应的 Jordan 块为
	\[
	J_{s_1}(\lambda),J_{s_2}(\lambda),\ldots,J_{s_m}(\lambda),
	\]
	则
	\[
	\nu_k=\sum_{i=1}^m \min\{k,s_i\}.
	\]
	因此
$\nu_k-\nu_{k-1}$ 等于大小至少为 $k$ 的 Jordan 块个数.
	
\end{proposition}

\begin{proof}
	由于相似不改变零空间维数，可设 $A$ 本身已经是 Jordan 标准型.先考虑一个 Jordan 块
	\[
	J_s(\lambda)=
	\begin{pmatrix}
		\lambda&0&0&\cdots&0\\
		1&\lambda&0&\cdots&0\\
		0&1&\lambda&\cdots&0\\
		\vdots&\vdots&\ddots&\ddots&\vdots\\
		0&0&\cdots&1&\lambda
	\end{pmatrix}.
	\]
	令
	\[
	N_s=J_s(\lambda)-\lambda E_s.
	\]
	则
	\[
	N_s=
	\begin{pmatrix}
		0&0&0&\cdots&0\\
		1&0&0&\cdots&0\\
		0&1&0&\cdots&0\\
		\vdots&\vdots&\ddots&\ddots&\vdots\\
		0&0&\cdots&1&0
	\end{pmatrix}.
	\]
	取标准基 $e_1,\ldots,e_s$，则
	\[
	N_s e_i=e_{i+1}\quad(1\leq i<s),\qquad N_s e_s=0.
	\]
	于是
	\[
	N_s^k e_i=
	\begin{cases}
		e_{i+k},& i+k\leq s,\\
		0,& i+k>s.
	\end{cases}
	\]
	所以当 $k\leq s$ 时，
	\[
	\ker N_s^k=\langle e_{s-k+1},\ldots,e_s\rangle,
	\]
	从而
	\[
	\dim\ker N_s^k=k.
	\]
	当 $k>s$ 时，$N_s^k=0$，故
	\[
	\dim\ker N_s^k=s.
	\]
	综上，
	\[
	\dim\ker\bigl(J_s(\lambda)-\lambda E_s\bigr)^k
	=
	\min\{k,s\}.
	\]
	
	现在设 $A$ 关于 $\lambda$ 的 Jordan 块为
	\[
	J_{s_1}(\lambda),J_{s_2}(\lambda),\ldots,J_{s_m}(\lambda).
	\]
	由于 Jordan 标准型是准对角阵，$(A-\lambda E)^k$ 仍是准对角阵.因此其零空间是各个对角块零空间的直和，所以维数相加：
	\[
	\dim\ker(A-\lambda E)^k
	=
	\sum_{i=1}^m
	\dim\ker\bigl(J_{s_i}(\lambda)-\lambda E_{s_i}\bigr)^k.
	\]
	于是
	\[
	\nu_k=\sum_{i=1}^m\min\{k,s_i\}.
	\]
	注意其他特征值 $\mu\neq\lambda$ 的 Jordan 块不贡献零空间，因为
	\[
	J_t(\mu)-\lambda E_t=(\mu-\lambda)E_t+N_t
	\]
	是可逆矩阵.
	
	最后考察差分.对一个大小为 $s$ 的 Jordan 块，有
	\[
	\min\{k,s\}-\min\{k-1,s\}
	=
	\begin{cases}
		1,&s\geq k,\\
		0,&s<k.
	\end{cases}
	\]
	因此，每个大小至少为 $k$ 的 Jordan 块贡献 $1$，每个大小小于 $k$ 的 Jordan 块贡献 $0$.把所有 Jordan 块的贡献相加，得到
$\nu_k-\nu_{k-1}$ 等于大小至少为 $k$ 的 Jordan 块个数.
	证毕.
	
\end{proof}

\begin{example}
	设 $T\in M_5(\F)$ 只有一个特征值 $3$.记
	\[
	N=T-3E.
	\]
	若
	\[
	\dim\ker N=2,\qquad
	\dim\ker N^2=3,\qquad
	\dim\ker N^3=4,\qquad
	\dim\ker N^4=5,
	\]
	则
	\[
	\nu_1=2,\qquad \nu_2=3,\qquad \nu_3=4,\qquad \nu_4=5.
	\]
	计算差分得
	\[
	\nu_1-\nu_0=2,\qquad
	\nu_2-\nu_1=1,\qquad
	\nu_3-\nu_2=1,\qquad
	\nu_4-\nu_3=1.
	\]
	因此大小至少为 $1$ 的 Jordan 块有 $2$ 个，大小至少为 $2$ 的 Jordan 块有 $1$ 个，大小至少为 $3$ 的 Jordan 块有 $1$ 个，大小至少为 $4$ 的 Jordan 块有 $1$ 个.
	
	所以 $T$ 关于特征值 $3$ 的 Jordan 块大小为
	\[
	4,\ 1.
	\]
	故
	\[
	T\sim \operatorname{diag}\bigl(J_4(3),J_1(3)\bigr).
	\]
	
\end{example}

\begin{example}
	设 $S\in M_5(\F)$ 有两个特征值 $2,-1$.对特征值 $2$，已知
	\[
	\dim\ker(S-2E)=2,\qquad
	\dim\ker(S-2E)^2=4.
	\]
	于是
	\[
	\nu_1-\nu_0=2,\qquad
	\nu_2-\nu_1=2.
	\]
	所以特征值 $2$ 对应两个 Jordan 块，并且这两个块的大小都至少为 $2$.
	
	又对特征值 $-1$，有
	\[
	\dim\ker(S+E)=1,
	\]
	故 $-1$ 至少对应一个 Jordan 块.由于 $S$ 是 $5$ 阶矩阵，总维数为 $5$，所以特征值 $2$ 只能占 $4$ 维，特征值 $-1$ 占 $1$ 维.因此 $2$ 对应的两个 Jordan 块大小为
	\[
	2,\ 2,
	\]
	而 $-1$ 对应的 Jordan 块大小为
	\[
	1.
	\]
	故
	\[
	S\sim \operatorname{diag}\bigl(J_2(2),J_2(2),J_1(-1)\bigr).
	\]
	
\end{example}


\section{一些有价值的习题}

\begin{problem}[11]
	设 $A$ 是一个 $n\times n$ 矩阵，它在 $n$ 维线性空间 $V$ 的一组基下定义线性变换.证明
	\[
	\dim(\Im A\cap\ker A)=\dim(\Im A)-\dim(\Im A^2).
	\]
	并推出：
	\begin{enumerate}
		\item $\ker A\subset \Im A$ 当且仅当
		\[
		\dim(\ker A)=\dim(\Im A)-\dim(\Im A^2);
		\]
		\item $\Im A\subset \ker A$ 当且仅当 $A^2=0$；
		\item $\Im A=\ker A$ 当且仅当 $A^2=0$ 且 $\dim(\ker A)=\dim(\Im A)$；
		\item $\dim(\Im A\cap\ker A)=0$ 当且仅当
		\[
		\dim(\Im A)=\dim(\Im A^2);
		\]
		\item $V=\Im A\oplus\ker A$ 当且仅当
		\[
		\dim(\Im A)=\dim(\Im A^2).
		\]
	\end{enumerate}
\end{problem}

\begin{proof}
	考虑限制映射
	\[
	A|_{\Im A}:\Im A\longrightarrow \Im A.
	\]
	它的核与像分别为
	\[
	\ker(A|_{\Im A})=\Im A\cap\ker A,
	\qquad
	\Im(A|_{\Im A})=\Im A^2.
	\]
	由维数公式得
	\[
	\dim(\Im A\cap\ker A)+\dim(\Im A^2)=\dim(\Im A),
	\]
	即
	\[
	\dim(\Im A\cap\ker A)=\dim(\Im A)-\dim(\Im A^2).
	\]
	下面证明各个结论.
	\begin{enumerate}
		\item 因为 $\ker A\subset \Im A$ 当且仅当
		\[
		\ker A=\Im A\cap\ker A,
		\]
		所以由上式立得结论.
		\item 有
		\[
		\Im A\subset \ker A
		\Longleftrightarrow
		A(\Im A)=0
		\Longleftrightarrow
		\Im A^2=0
		\Longleftrightarrow A^2=0.
		\]
		\item 由 (2)，$\Im A\subset\ker A$ 当且仅当 $A^2=0$.此时再由二者维数相等，即得
		$\Im A=\ker A$.
		\item 由已证等式直接得到.
		\item 由维数公式
		\[
		\dim V=\dim(\Im A)+\dim(\ker A),
		\]
		可知 $V=\Im A\oplus\ker A$ 当且仅当
		$\Im A\cap\ker A=0$.再由 (4) 得证.
	\end{enumerate}
\end{proof}

\begin{problem}[12]
	若两个实方阵在复数域上相似，则它们在实数域上相似.
\end{problem}

\begin{proof}
	设 $A,B\in M_n(\mathbb{R})$，且存在 $P\in \GL_n(\mathbb{C})$ 使得
	\[
	A=PBP^{-1},
	\]
	即
	\[
	AP=PB.
	\]
	写
	\[
	P=Q+iR,
	\qquad Q,R\in M_n(\mathbb{R}).
	\]
	则
	\[
	A(Q+iR)=(Q+iR)B,
	\]
	比较实部与虚部，得
	\[
	AQ=QB,
	\qquad
	AR=RB.
	\]
	考虑实系数多项式
	\[
	p(t)=\det(Q+tR).
	\]
	若 $p(t)$ 恒为零，则 $p(i)=\det(Q+iR)=\det P=0$，矛盾.故可取 $t_0\in\mathbb{R}$ 使得
	\[
	\det(Q+t_0R)\neq0.
	\]
	令 $C=Q+t_0R$，则 $C\in \GL_n(\mathbb{R})$，且
	\[
	AC=CB.
	\]
	所以
	\[
	A=CBC^{-1}.
	\]
	故 $A,B$ 在实数域上相似.
\end{proof}

\begin{problem}[15]
	设 $A$ 是一个数域上的 $n\times n$ 矩阵.证明
	\[
	\deg(\lambda E_n-A)^*=n-1.
	\]
\end{problem}

\begin{proof}
	由伴随矩阵的定义，
	\[
	(\lambda E_n-A)(\lambda E_n-A)^*=|\lambda E_n-A|E_n.
	\]
	设
	\[
	(\lambda E_n-A)^*=P_m\lambda^m+P_{m-1}\lambda^{m-1}+\cdots+P_0,
	\qquad P_m\neq0.
	\]
	又
	\[
	|\lambda E_n-A|=\lambda^n+a_{n-1}\lambda^{n-1}+\cdots+a_0.
	\]
	于是左边最高次项为 $P_m\lambda^{m+1}$，右边最高次项为 $E_n\lambda^n$，比较次数得
	\[
	m+1=n.
	\]
	因此
	\[
	\deg(\lambda E_n-A)^*=n-1.
	\]
\end{proof}

\begin{problem}[16]
	设 $A(\lambda)$ 是一个 $\lambda$-矩阵.若 $\deg A(\lambda)=d$，是否必有
	\[
	\deg |A(\lambda)|\geq d?
	\]
\end{problem}

\begin{solution}
	不一定.例如取
	\[
	A(\lambda)=
	\begin{pmatrix}
		\lambda^2+1&\lambda^2\\
		\lambda^2&\lambda^2-1
	\end{pmatrix}.
	\]
	则
	\[
	\deg A(\lambda)=2,
	\]
	但
	\[
	|A(\lambda)|=(\lambda^2+1)(\lambda^2-1)-\lambda^4=-1,
	\]
	故
	\[
	\deg |A(\lambda)|=0<2.
	\]
\end{solution}

\begin{problem}[17]
	若 $n\times n$ 的 $\lambda$-矩阵 $A(\lambda)$ 可逆，则 $A(\lambda)$ 的秩为 $n$.反之是否成立？
\end{problem}

\begin{solution}
	若 $A(\lambda)$ 可逆，则
	\[
	\det A(\lambda)\neq0,
	\]
	从而 $A(\lambda)$ 有一个 $n$ 阶子式不为零，故
	\[
	r(A(\lambda))=n.
	\]
	反之不成立.例如取
	\[
	A(\lambda)=\lambda E_n-B,
	\qquad B\in M_n(\F).
	\]
	由于
	\[
	\det(\lambda E_n-B)
	\]
	是 $B$ 的特征多项式，故它是一个非零的 $n$ 次多项式，从而
	\[
	r(\lambda E_n-B)=n.
	\]
	但是 $\det(\lambda E_n-B)$ 不是非零常数，所以 $\lambda E_n-B$ 在 $\lambda$-矩阵意义下不可逆.
\end{solution}

\begin{problem}[34]
	设 $\lambda_1,\lambda_2,\ldots,\lambda_n$ 是 $n\times n$ 矩阵 $A$ 的全部特征值，重数也计入其中.证明：对任意多项式 $g(x)$，有
	\[
	|g(A)|=\prod_{i=1}^n g(\lambda_i).
	\]
\end{problem}

\begin{proof}
	在包含所有特征值的扩域中，$A$ 可三角化.设 $A$ 与上三角矩阵 $U$ 相似，且 $U$ 的对角元依次为
	\[
	\lambda_1,\lambda_2,\ldots,\lambda_n.
	\]
	于是 $g(A)$ 与 $g(U)$ 相似，而 $g(U)$ 仍是上三角矩阵，且其对角元为
	\[
	g(\lambda_1),g(\lambda_2),\ldots,g(\lambda_n).
	\]
	所以
	\[
	|g(A)|=\prod_{i=1}^n g(\lambda_i).
	\]
\end{proof}

\begin{problem}[45]
	设 $A\in M_n(\F)$，$m_A(\lambda)$ 为 $A$ 的极小多项式.若 $g(\lambda)\in\F[\lambda]$，证明：
	\[
	(g(\lambda),m_A(\lambda))=1
	\Longleftrightarrow
	\det g(A)\neq0.
	\]
\end{problem}

\begin{proof}
	若 $(g,m_A)=1$，则存在 $u(\lambda),v(\lambda)\in\F[\lambda]$，使得
	\[
	u(\lambda)g(\lambda)+v(\lambda)m_A(\lambda)=1.
	\]
	代入 $A$ 得
	\[
	u(A)g(A)+v(A)m_A(A)=E_n.
	\]
	由于 $m_A(A)=0$，故
	\[
	u(A)g(A)=E_n,
	\]
	所以 $g(A)$ 可逆，因而 $\det g(A)\neq0$.
	
	反过来，若 $\det g(A)\neq0$，设 $\lambda_1,\ldots,\lambda_n$ 是 $A$ 的特征值.由上一题，
	\[
	\det g(A)=\prod_{i=1}^n g(\lambda_i)\neq0,
	\]
	故 $g(\lambda_i)\neq0$ 对所有 $i$ 成立.而 $m_A$ 的根正是 $A$ 的特征值，故 $g$ 与 $m_A$ 无公共根，从而
	\[
	(g,m_A)=1.
	\]
\end{proof}

\begin{problem}[59]
	设
	\[
	L=L_1\oplus L_2,
	\]
	其中 $L_1,L_2$ 是 $A$-子空间.证明
	\[
	f_A(\lambda)=f_{A|_{L_1}}(\lambda)f_{A|_{L_2}}(\lambda),
	\]
	其中 $f_A(\lambda)$ 是 $A$ 的特征多项式.
\end{problem}

\begin{proof}
	设 $\dim L_1=r$，$\dim L=s$.取 $L_1$ 的一组基
	\[
	\xi_1,\ldots,\xi_r,
	\]
	以及 $L_2$ 的一组基
	\[
	\xi_{r+1},\ldots,\xi_s.
	\]
	由于 $L=L_1\oplus L_2$，故
	\[
	\xi_1,\ldots,\xi_r,\xi_{r+1},\ldots,\xi_s
	\]
	是 $L$ 的一组基.
	
	由于 $L_i$ 是 $A$-子空间，故
	\[
	A|_{L_i}\in\End(L_i),\qquad i=1,2.
	\]
	设 $A|_{L_1}$ 在基 $\xi_1,\ldots,\xi_r$ 下的矩阵为 $A_1$，$A|_{L_2}$ 在基 $\xi_{r+1},\ldots,\xi_s$ 下的矩阵为 $A_2$，即
	\[
	(A|_{L_1}\xi_1\ \cdots\ A|_{L_1}\xi_r)=(\xi_1\ \cdots\ \xi_r)A_1,
	\]
	\[
	(A|_{L_2}\xi_{r+1}\ \cdots\ A|_{L_2}\xi_s)=(\xi_{r+1}\ \cdots\ \xi_s)A_2.
	\]
	又对任意 $\alpha\in L_i$，有 $A|_{L_i}\alpha=A\alpha$.所以
	\[
	(A\xi_1\ \cdots\ A\xi_r\ A\xi_{r+1}\ \cdots\ A\xi_s)
	=(\xi_1\ \cdots\ \xi_r\ \xi_{r+1}\ \cdots\ \xi_s)
	\begin{pmatrix}
		A_1&0\\
		0&A_2
	\end{pmatrix}.
	\]
	即 $A$ 在这组基下的矩阵为
	\[
	\begin{pmatrix}
		A_1&0\\
		0&A_2
	\end{pmatrix}.
	\]
	因此
	\[
	\begin{aligned}
		f_A(\lambda)
		&=\left|\lambda E_s-
		\begin{pmatrix}
			A_1&0\\
			0&A_2
		\end{pmatrix}\right|\\
		&=\left|\begin{matrix}
			\lambda E_r-A_1&0\\
			0&\lambda E_{s-r}-A_2
		\end{matrix}\right|\\
		&=|\lambda E_r-A_1|\,|\lambda E_{s-r}-A_2|\\
		&=f_{A|_{L_1}}(\lambda)f_{A|_{L_2}}(\lambda).
	\end{aligned}
	\]
\end{proof}



\begin{problem}[65]
	若两个数值矩阵 $A,B$ 作为 $\lambda$-矩阵相抵，则它们作为数值矩阵也相抵.
\end{problem}

\begin{proof}
	法一：利用 $\lambda$-矩阵的初等变换不改变 $\lambda$-矩阵的秩.若 $A,B\in M_n(\F)$ 作为 $\lambda$-矩阵相抵，则它们作为 $\lambda$-矩阵的秩相等.由于 $A,B$ 是常数矩阵，它们作为 $\lambda$-矩阵时非零子式的最大阶数，正是作为数值矩阵时非零子式的最大阶数.故
	\[
	r(A)=r(B).
	\]
	由矩阵相抵的判别，$A$ 与 $B$ 作为数值矩阵相抵.
	
	法二：设存在可逆的 $\lambda$-矩阵 $P(\lambda),Q(\lambda)$，使得
	\[
	A=P(\lambda)BQ(\lambda).
	\]
	写
	\[
	P(\lambda)=\lambda^rP_r+\cdots+\lambda P_1+P_0,
	\qquad
	Q(\lambda)=\lambda^sQ_s+\cdots+\lambda Q_1+Q_0.
	\]
	由下面的引理可知 $P_0,Q_0$ 可逆.比较等式两边的常数项，得
	\[
	A=P_0BQ_0.
	\]
	故 $A,B$ 作为数值矩阵相抵.
\end{proof}

\begin{lemma}
	设
	\[
	M(\lambda)=\lambda^mM_m+\cdots+\lambda M_1+M_0
	\]
	是 $n$ 阶可逆 $\lambda$-矩阵，其中 $M_i$ 是数值矩阵，则 $M_0$ 可逆.
\end{lemma}

\begin{proof}
	设
	\[
	M(\lambda)^{-1}=\lambda^uN_u+\cdots+\lambda N_1+N_0.
	\]
	代入
	\[
	M(\lambda)M(\lambda)^{-1}=E_n
	\]
	并比较常数项，得到
	\[
	M_0N_0=E_n.
	\]
	故 $M_0$ 可逆.
\end{proof}

\begin{problem}[85]
	设 $A$ 是一个 $4\times4$ 矩阵，特征多项式为 $(\lambda-1)^4$.
	\begin{enumerate}
		\item 证明 $A$ 的 Jordan 标准型有五种可能，但极小多项式只有四种可能.
		\item 证明 $A$ 的 Jordan 标准型由极小多项式的次数与 $r(E-A)$ 共同决定.
	\end{enumerate}
\end{problem}

\begin{proof}
	\begin{enumerate}
		\item 由于 $\mathcal{P}(4)=5$，即
		\[
		4=4=3+1=2+2=2+1+1=1+1+1+1,
		\]
		所以 Jordan 标准型有五种可能：
		\[
		J_4(1),
		\]
		\[
		\operatorname{diag}(J_3(1),J_1(1)),
		\]
		\[
		\operatorname{diag}(J_2(1),J_2(1)),
		\]
		\[
		\operatorname{diag}(J_2(1),J_1(1),J_1(1)),
		\]
		\[
		\operatorname{diag}(J_1(1),J_1(1),J_1(1),J_1(1)).
		\]
		它们对应的极小多项式分别为
		\[
		(\lambda-1)^4,
		\quad
		(\lambda-1)^3,
		\quad
		(\lambda-1)^2,
		\quad
		(\lambda-1)^2,
		\quad
		(\lambda-1).
		\]
		因此 Jordan 标准型有五种，而极小多项式只有四种.
		
		\item 一般地，特征多项式与极小多项式不能完全决定 Jordan 标准型.在本题中，只有当极小多项式为 $(\lambda-1)^2$ 时，需要区分两种情形：
		\[
		\operatorname{diag}(J_2(1),J_2(1))
		\quad\text{与}\quad
		\operatorname{diag}(J_2(1),J_1(1),J_1(1)).
		\]
		此时
		\[
		r(E-\operatorname{diag}(J_2(1),J_2(1)))=2,
		\]
		而
		\[
		r(E-\operatorname{diag}(J_2(1),J_1(1),J_1(1)))=1.
		\]
		其余极小多项式次数对应的 Jordan 标准型唯一.因此 $A$ 的 Jordan 标准型由极小多项式的次数与 $r(E-A)$ 共同决定.
	\end{enumerate}
\end{proof}

\begin{problem}[88]
	设 $\F$ 是一个数域.对任意 $a_0,a_1,\ldots,a_{n-1}\in\F$，定义 $n\times n$ 循环矩阵
	\[
	A=C[a_0,a_1,\ldots,a_{n-1}]
	=
	\begin{pmatrix}
		a_0&a_1&a_2&\cdots&a_{n-1}\\
		a_{n-1}&a_0&a_1&\cdots&a_{n-2}\\
		a_{n-2}&a_{n-1}&a_0&\cdots&a_{n-3}\\
		\vdots&\vdots&\vdots&&\vdots\\
		a_1&a_2&a_3&\cdots&a_0
	\end{pmatrix}.
	\]
	证明：$A$ 可逆当且仅当多项式
	\[
	f(x)=\sum_{i=0}^{n-1}a_ix^i
	\]
	与 $x^n-1$ 互素.
\end{problem}

\begin{proof}
	设
	\[
	f(x)=\sum_{i=0}^{n-1}a_ix^i.
	\]
	在包含 $x^n-1$ 全部根的扩域中，记
	\[
	\varepsilon_0=1,\varepsilon_1,\ldots,\varepsilon_{n-1}
	\]
	为 $x^n-1$ 的全部根.由于数域的特征为 $0$，这些根互不相同.
	
	取 Vandermonde 矩阵
	\[
	P=
	\begin{pmatrix}
		1&1&\cdots&1\\
		1&\varepsilon_1&\cdots&\varepsilon_{n-1}\\
		\vdots&\vdots&&\vdots\\
		1&\varepsilon_1^{n-1}&\cdots&\varepsilon_{n-1}^{n-1}
	\end{pmatrix}.
	\]
	由 Vandermonde 行列式知 $|P|\neq0$.直接计算可得
	\[
	AP=P
	\begin{pmatrix}
		f(1)&&&\\
		&f(\varepsilon_1)&&\\
		&&\ddots&\\
		&&&f(\varepsilon_{n-1})
	\end{pmatrix}.
	\]
	于是
	\[
	|A|=f(1)f(\varepsilon_1)\cdots f(\varepsilon_{n-1}).
	\]
	因此 $A$ 可逆当且仅当
	\[
	f(\varepsilon_i)\neq0,
	\qquad i=0,1,\ldots,n-1.
	\]
	这等价于 $f(x)$ 与 $x^n-1$ 没有公共根，即
	\[
	(f(x),x^n-1)=1.
	\]
\end{proof}

\begin{problem}[89]
	证明：线性变换 $A$ 关于特征值 $\lambda_0$ 的 Jordan 块个数，等于属于 $\lambda_0$ 的特征子空间的维数.
\end{problem}

\begin{proof}
	设 $A$ 的 Jordan 标准型中，属于 $\lambda_0$ 的 Jordan 块为
	\[
	J_{r_1}(\lambda_0),\ldots,J_{r_t}(\lambda_0).
	\]
	其他特征值对应的 Jordan 块合在一起记为 $J'$.则
	\[
	A\sim \operatorname{diag}(J_{r_1}(\lambda_0),\ldots,J_{r_t}(\lambda_0),J').
	\]
	对每个块有
	\[
	r(\lambda_0E_{r_i}-J_{r_i}(\lambda_0))=r_i-1.
	\]
	而 $\lambda_0E-J'$ 中的每个对角块均可逆.因此
	\[
	r(\lambda_0E-A)=n-t.
	\]
	于是
	\[
	\dim\ker(\lambda_0E-A)=n-r(\lambda_0E-A)=t.
	\]
	左边正是属于 $\lambda_0$ 的特征子空间维数，右边正是属于 $\lambda_0$ 的 Jordan 块个数.
\end{proof}

\begin{example}[Jordan 基的一种求法]
	给定矩阵
	\[
	A=
	\begin{pmatrix}
		8&-3&3&-2\\
		0&3&0&2\\
		-10&6&-3&4\\
		2&0&1&2
	\end{pmatrix}.
	\]
	可算得 $A$ 的特征多项式为
	\[
	f_A(\lambda)=(\lambda-2)^2(\lambda-3)^2.
	\]
	对 $\lambda E-A$ 作初等变换，可得 $A$ 的初等因子为
	\[
	(\lambda-2)^2,
	\qquad
	(\lambda-3)^2.
	\]
	因此 $A$ 的 Jordan 标准型为
	\[
	J=
	\begin{pmatrix}
		2&0&0&0\\
		1&2&0&0\\
		0&0&3&0\\
		0&0&1&3
	\end{pmatrix}.
	\]
	
	若 $P=(\xi_1\ \xi_2\ \xi_3\ \xi_4)$，且 $AP=PJ$，则
	\[
	(A\xi_1\ A\xi_2\ A\xi_3\ A\xi_4)
	=(\xi_1\ \xi_2\ \xi_3\ \xi_4)
	\begin{pmatrix}
		2&0&0&0\\
		1&2&0&0\\
		0&0&3&0\\
		0&0&1&3
	\end{pmatrix},
	\]
	即
	\[
	\left\{
	\begin{aligned}
		A\xi_1&=2\xi_1+\xi_2,\\
		A\xi_2&=2\xi_2,\\
		A\xi_3&=3\xi_3+\xi_4,\\
		A\xi_4&=3\xi_4.
	\end{aligned}
	\right.
	\]
	这说明 $\xi_2,\xi_4$ 分别是属于特征值 $2,3$ 的特征向量.
	
	对 $\lambda_1=2$，解
	\[
	(2E-A)X=0
	\]
	可取
	\[
	\eta_1=
	\begin{pmatrix}
		1\\0\\-2\\0
	\end{pmatrix}.
	\]
	再解
	\[
	(A-2E)X=\eta_1,
	\]
	其通解为
	\[
	X=\frac14
	\begin{pmatrix}
		0\\-2\\0\\1
	\end{pmatrix}
	+t
	\begin{pmatrix}
		1\\0\\-2\\0
	\end{pmatrix},
	\qquad t\in\mathbb{C}.
	\]
	取 $t=0$，得
	\[
	\eta_2=
	\begin{pmatrix}
		0\\-rac12\\0\\\frac14
	\end{pmatrix}.
	\]
	
	对 $\lambda_2=3$，解
	\[
	(3E-A)X=0
	\]
	可取
	\[
	\zeta_1=
	\begin{pmatrix}
		-3\\1\\6\\0
	\end{pmatrix}.
	\]
	再解
	\[
	(A-3E)X=\zeta_1,
	\]
	其通解为
	\[
	X=\frac16
	\begin{pmatrix}
		0\\7\\3\\3
	\end{pmatrix}
	+s
	\begin{pmatrix}
		-3\\1\\6\\0
	\end{pmatrix},
	\qquad s\in\mathbb{C}.
	\]
	取 $s=0$，得
	\[
	\zeta_2=\frac16
	\begin{pmatrix}
		0\\7\\3\\3
	\end{pmatrix}.
	\]
	于是可取
	\[
	P=(\eta_2\ \eta_1\ \zeta_2\ \zeta_1)
	=
	\begin{pmatrix}
		0&1&0&-3\\
		-\frac12&0&\frac76&1\\
		0&-2&\frac12&6\\
		\frac14&0&\frac12&0
	\end{pmatrix}.
	\]
	则
	\[
	P^{-1}AP=J.
	\]
\end{example}

\begin{proposition}
	设 $A\in M_n(\F)$.则 $A$ 在 $\F$ 上可对角化当且仅当下列条件之一成立：
	\begin{enumerate}
		\item $A$ 在 $\F^n$ 中有 $n$ 个线性无关的特征向量；
		\item $A$ 的极小多项式在 $\F$ 上可以分解为互不相同的一次因式的乘积；
		\item $A$ 的初等因子全是 $\F$ 上的一次因式；
		\item $\F^n$ 可以分解为 $A$ 的特征子空间的直和；
		\item $A$ 的每个特征值都属于 $\F$，且每个特征值的几何重数等于代数重数.
	\end{enumerate}
\end{proposition}
